\documentclass[11pt]{article}
\usepackage[margin=2.3cm]{geometry}

\usepackage{graphicx}
\graphicspath{{../diagrams/}{diagrams/}}
\usepackage{float}
\usepackage{amsmath,amssymb,amsfonts}
\usepackage{hyperref}
\usepackage{booktabs}
\usepackage{xcolor}
\usepackage{array}
\usepackage{bm}

\hypersetup{
  colorlinks=true,
  linkcolor=blue!50!black,
  citecolor=green!50!black,
  urlcolor=blue!50!black,
}

\newcommand{\su}{\mathfrak{su}}
\newcommand{\so}{\mathfrak{so}}
\newcommand{\Spin}{\mathrm{Spin}}
\newcommand{\half}{\tfrac{1}{2}}
\newcommand{\Adj}{\mathrm{Adj}}
\newcommand{\Sym}{\mathrm{Sym}}
\newcommand{\dimq}{\dim_q}
\newcommand{\ket}[1]{\left|#1\right\rangle}
\newcommand{\bra}[1]{\left\langle#1\right|}
\newcommand{\ip}[2]{\left\langle#1,#2\right\rangle}
\newcommand{\mPl}{m_{\mathrm{Pl}}}
\newcommand{\me}{m_e}
\newcommand{\hv}{h^{\vee}}
\newcommand{\rank}{\mathrm{rank}}
\newcommand{\Tr}{\mathrm{Tr}}
\newcommand{\C}{\mathbb{C}}
\newcommand{\Z}{\mathbb{Z}}
\newcommand{\F}{\mathbb{F}}
\newcommand{\R}{\mathbb{R}}
\newcommand{\PSL}{\mathrm{PSL}}
\newcommand{\Aut}{\mathrm{Aut}}
\newcommand{\SL}{\mathrm{SL}}
\newcommand{\GL}{\mathrm{GL}}
\newcommand{\SU}{\mathrm{SU}}
\newcommand{\U}{\mathrm{U}}
\newcommand{\SO}{\mathrm{SO}}
\newcommand{\GeV}{\mathrm{GeV}}
\newcommand{\MeV}{\mathrm{MeV}}

\begin{document}

\title{Sakharov-Induced Einstein Gravity\\[2pt]
from the Null Worldtube Condensate\\[4pt]
\normalsize (Paper 18)}

\author{
  James~P.~Galasyn$^{1}$ \quad and \quad Claude~Th\'eodore$^{2}$\\[2pt]
  \small $^1$Independent researcher \\
  \small $^2$Anthropic
}

\date{\today}

\maketitle

\begin{abstract}
We derive Einstein's equations as the long-wavelength dynamics
induced by the Null Worldtube Theory (NWT) condensate.  Starting
from the minimal three-field Lagrangian of Papers~15--16%
~\cite{paper15, paper16}---a complex scalar coupled to a $\U(1)$
gauge field and an $S^2$-valued Skyrme--Faddeev field at the BPS
critical coupling---we show that integrating out matter at one
loop on an arbitrary curved background generates the
Einstein--Hilbert action plus controlled higher-curvature
corrections, in the manner of Sakharov's induced
gravity~\cite{sakharov1968, visser2002}.

The graviton is \emph{not} a member of the topological vortex zoo:
the $\mathrm{L}1$ field content of the substrate carries no
spacetime spin-$2$ representation at linear order.  Rather, the
metric perturbation $h_{\mu\nu}$ is an emergent variable whose
kinetic action is generated by the matter loop, with the
Einstein--Hilbert coefficient $1/(16\pi G)$ set by the heat-kernel
$a_2$ coefficient of the matter spectrum.  At flat-space leading
order the induced gravitational coupling is
\begin{equation*}
  \frac{1}{16\pi G}
  = \frac{N_{\mathrm{DOF}}\,\Lambda_{\mathrm{UV}}^2}{192\pi^2},
\end{equation*}
with $N_{\mathrm{DOF}} = 6$ for the minimal matter content and
$\Lambda_{\mathrm{UV}} \approx \mPl$.

Combining the Sakharov coefficient with Paper~17's closed form
$\mPl = \me / [(8/7)(1 + \alpha/7 + 3\alpha^2)\,\alpha^{21/2}]$
gives the structural relation
$\Lambda_{\mathrm{UV}}/\mPl =
\sqrt{12\pi/N_{\mathrm{DOF}}^{\mathrm{sub}}}$, with
$N_{\mathrm{DOF}}^{\mathrm{sub}}$ a $c_J$-weighted sum over
\emph{substrate} matter fields (the bosonic three-field Lagrangian
of Paper~16, optionally extended to the $\so(10)$ GUT spectrum of
Paper~14~\cite{paper14}).  IR-projected SM fermions and massless
gauge bosons are topological projections of the substrate, not new
fundamental fields, and do not enter the heat-kernel sum directly;\
naively summing them with standard Birrell--Davies coefficients
reproduces the well-known Sakharov fermion-sign issue but does
not apply at the substrate level.  Identifying the substrate count
with one mode per $\so(10)/\mathrm{SM}$ generator gives
$N_{\mathrm{DOF}}^{\mathrm{sub}} = 6 + 33 = 39$ and
$\Lambda_{\mathrm{UV}} = 0.983\,\mPl$, consistent with
$\Lambda_{\mathrm{UV}} \approx \mPl$ to $1.7\%$.  A
first-principles derivation of the substrate $c_J$ is logged as a
follow-up.  The Newton's-constant prediction following from
Paper~17's closed form matches CODATA to $+29$\,ppm, inside the
$\pm 22$\,ppm experimental uncertainty on $G$.

The Paper~17 identity $\mPl^2/\me^2 = |A_{K_7}|^{-2}$ is identified
here as the structural relation between the ultraviolet cutoff and
the infrared electron sector.  The $10^{45}$ Phase~5 mismatch
reported in Paper~16---between the naive heat-kernel computation on
$S^3/2I$ and the experimentally required $1/G$---is structurally
identified with the missing $K_7$ Lie-theoretic factor that the
bare $S^3/2I$ Casimir computation did not include;\ a corrected
Phase~5 calculation incorporating this factor is sketched in
Section~\ref{sec:G3-phase5-sketch} and the full evaluation is
logged as a follow-up.

Variation of the matter-loop-induced effective action $\Gamma[g]$
yields the full nonlinear Einstein equations with cosmological
constant and Stelle higher-curvature corrections,
\begin{equation*}
  G_{\mu\nu} + \Lambda_{\mathrm{cc}}\,g_{\mu\nu}
  = 8\pi G\,T_{\mu\nu}
  + \alpha_S\,\partial_\mu\partial_\nu R + \cdots,
\end{equation*}
with $T_{\mu\nu}$ the matter stress-energy of the
$\mathrm{L}1{+}\mathrm{L}2{+}\mathrm{L}3$ sector and $G$ matching
Paper~17's structurally derived value.  Bianchi's identity
automatically enforces $\nabla^\mu T_{\mu\nu} = 0$.  The
Schwarzschild metric is verified as a vacuum solution by direct
symbolic computation: all $10$ independent components of the
symmetric Ricci tensor $R_{\mu\nu}$ vanish identically, evaluated
in the companion library's gravity shim.
Linearisation around flat space reproduces the gravitational-wave
equation $\Box\,h^{\mathrm{TT}}_{\mu\nu} = -16\pi G\,T^{\mathrm{TT}}_{\mu\nu}$,
with two physical transverse-traceless polarizations and Newton's
law in the non-relativistic limit.

We discuss known open problems:\ the cosmological-constant
fine-tuning ($\sim 10^{120}$ gap, of which the $K_7$-amplitude
factor closes $\sim 45$ orders, leaving $\sim 75$ unresolved and
deferred to Paper~21~\cite{paper19}); the formation of black holes
from dynamical collapse of NWT vortex configurations;\ ghost-freedom
of the Stelle higher-curvature sector beyond linear order;\ and a
quantum-level treatment of diffeomorphism invariance on curved
backgrounds.  None of these are NWT-specific; all are flagged as
candidate research programs for follow-up papers.

The result strengthens the framework of Papers~14--17: the
gravitational coupling that emerges from the NWT condensate at the
classical level is verified at the quantum-loop level (Sakharov),
at the structural level ($K_7$ amplitude), and now at the
dynamical level (Einstein equations on arbitrary backgrounds).
The non-linearity of gravity is geometric, following from
variation of the Sakharov-induced Einstein--Hilbert action, not
from additional matter physics.
\end{abstract}

\tableofcontents

% ==========================================================
\section{Introduction}
\label{sec:intro}
% ==========================================================

\subsection{Context: the NWT gravitational program}

Null Worldtube Theory~\cite{paper14, paper15, paper16, paper17}
is a topological-substrate reformulation of the Standard Model and
gravity in which particles are torus-knot vortex defects of a
relativistic condensate $\psi$, and the substrate algebra
($\mathrm{Cl}(0,7)$ octonion Clifford,
$\mathrm{Spin}(7)$ holonomy, the binary icosahedral group $2I$, and
the complete graph $K_7$) generates Standard-Model and gravitational
observables structurally.  The gravitational program of NWT has
proceeded across four papers, each tightening the predicted value
of Newton's constant $G$.

Paper~14~\cite{paper14} established the leading-order structural
form
\begin{equation}
  G_{\mathrm{LO}} \;=\; (8/7)^2\,\alpha^{21}\,\frac{\hbar c}{\me^2},
  \label{eq:p14-LO}
\end{equation}
accurate to $0.24\%$ of CODATA, identifying $\alpha^{21}$ as the
heptafoil Wilson amplitude on the $K_7$ graph and $(8/7)$ as the
$\mathrm{Spin}(7)$ spinor-to-vector dimension ratio.
Paper~15~\cite{paper15} placed this expression on the
Brieskorn--Poincar\'e sphere $S^3/2I$ as a one-loop Casimir
computation, identifying the BPS-trefoil background that supports
the calculation.  Paper~16~\cite{paper16} supplied the matter
Lagrangian on which the loop is taken---a complex scalar $\psi$
coupled to a $\U(1)$ gauge field $A_\mu$ and an $S^2$-valued
Skyrme--Faddeev field $n^a$ at the BPS critical coupling
$\lambda = e^2/2$---and reported a $10^{45}$ structural gap
between the naive heat-kernel computation on $S^3/2I$ and the
experimentally required $1/G$, deferring its resolution.
Paper~17~\cite{paper17} closed the structural derivation at
next-to-next-to-leading order via $K_7$ graph-state information
theory:
\begin{equation}
  \frac{\me}{\mPl}
    \;=\;
    \frac{\dim S}{\dim V}\,
    \left[\, 1
      + \frac{\alpha}{\dim V}
      + \frac{\dim \Adj}{\dim S}\,\alpha^2
    \right]
    \sin^{21}\!\left(\frac{\pi}{k + \hv}\right),
  \label{eq:p17}
\end{equation}
where $V$, $S$, $\Adj$ are the $7$-dimensional vector,
$8$-dimensional spinor, and $21$-dimensional adjoint representations
of $\mathrm{Spin}(7)$, and $\hv = 5$ is its dual Coxeter number.
Equation~\eqref{eq:p17} agrees with CODATA $\me/\mPl$ to $-0.0001\%$
and predicts $G$ within $11$\,ppm of CODATA, inside the experimental
error bar on $G$ itself.

What Papers 14--17 deliver is a \emph{number}.  They produce the
structurally-correct value of $G$ from substrate primitives, but
they do not produce the dynamics in which that number lives.
$G$ is the coefficient of the Einstein--Hilbert action; that action
has to be derived, not asserted.  The present paper supplies the
dynamical content: integrating out the matter sector at one loop on
an arbitrary curved background generates the Einstein--Hilbert
action plus controlled higher-curvature corrections, with the
coefficient $1/(16\pi G)$ consistent with Paper~17's structural
prediction at leading order.

\subsection{The Sakharov idea}

The mechanism is Andrei Sakharov's 1968 proposal that gravity is
not a fundamental interaction but the induced low-energy dynamics
of any theory whose matter sector is coupled to a background
metric~\cite{sakharov1968}.  In Sakharov's analysis, the one-loop
effective action obtained by integrating out matter on an
arbitrary curved $g_{\mu\nu}$ contains, in its derivative
expansion, a term proportional to the Ricci scalar:
\begin{equation}
  \Gamma[g]_{\mathrm{matter\;loop}}
   \;\supset\;
   \frac{1}{16\pi G_{\mathrm{ind}}}
     \int d^4 x\,\sqrt{-g}\,R(g) \;+\; \cdots,
\end{equation}
with the induced Newton constant $G_{\mathrm{ind}}$ fixed by the
heat-kernel $a_2$ coefficient of the matter spectrum.  At
flat-space leading order,
\begin{equation}
  \frac{1}{16\pi G_{\mathrm{ind}}}
    \;=\;
    \frac{N_{\mathrm{DOF}}\,\Lambda_{\mathrm{UV}}^2}{192\pi^2},
  \label{eq:sakharov-LO}
\end{equation}
where $N_{\mathrm{DOF}}$ counts the matter degrees of freedom
participating in the loop and $\Lambda_{\mathrm{UV}}$ is the cutoff.
Subsequent work has placed the construction on a firm modern
footing (Visser 2002~\cite{visser2002}; Stelle's higher-curvature
classification 1977~\cite{stelle1977}; Vassilevich's heat-kernel
manual 2003~\cite{vassilevich2003}), and Volovik's superfluid-helium
program~\cite{volovik2003} has demonstrated the same mechanism in
laboratory analog systems.

NWT belongs to this lineage but with a specific structural
specification: the matter sector is the BPS three-field theory of
Paper~16, and the cutoff $\Lambda_{\mathrm{UV}}$ is fixed by
Paper~17's $K_7$-amplitude prediction
\eqref{eq:p17} rather than left as a free parameter.  The induced
$G$ is therefore not just calculable in form but pinned in
magnitude.

\subsection{What this paper proves}

We establish the following, in order:

\begin{enumerate}
\item The covariant matter action $S_{\mathrm{matter}}[\psi, A, n; g]$
      of Paper~16 has a unique homogeneous BPS vacuum on
      $g_{\mu\nu} = \eta_{\mu\nu}$, with vanishing matter
      stress-energy $T^{(M)}_{\mu\nu} = 0$ at the vacuum point.
      Linearisation $g = \eta + h$ produces, to first order in
      $h$, a Sakharov-type matter--graviton coupling
      $-\tfrac{1}{2}\,h^{\mu\nu}\,T^{(M)}_{\mu\nu}$
      where $T^{(M)}_{\mu\nu}$ is computed at second order in
      matter fluctuations around the BPS vacuum.
\item Integrating out matter at one loop on $g = \eta + h$
      generates the linearized graviton kinetic term, and the
      coefficient $1/(16\pi G)$ is fixed by the heat-kernel $a_2$
      coefficient of the matter spectrum.
\item Equating the Sakharov coefficient with Paper~17's closed
      form for $\me/\mPl$ gives the structural relation
      $\Lambda_{\mathrm{UV}}/\mPl =
      \sqrt{12\pi/N_{\mathrm{DOF}}^{\mathrm{sub}}}$, with
      $N_{\mathrm{DOF}}^{\mathrm{sub}}$ summed over the
      \emph{substrate} matter fields (the bosonic three-field
      Lagrangian of Paper~16, optionally extended to the
      $\so(10)$ GUT spectrum of Paper~14~\cite{paper14}).  IR-projected
      SM fermions and massless gauge bosons are topological
      projections of the substrate, not new fundamental fields;\
      summing them with standard Birrell--Davies coefficients
      reproduces the well-known Sakharov fermion-sign
      issue~\cite{adler1982, visser2002} but does not apply at the
      substrate level.  Identifying the substrate count with one
      mode per $\so(10)/\mathrm{SM}$ generator gives
      $N_{\mathrm{DOF}}^{\mathrm{sub}} = 6 + 33 = 39$ and
      $\Lambda_{\mathrm{UV}} = 0.983\,\mPl$, consistent with
      $\Lambda_{\mathrm{UV}} \approx \mPl$ to $1.7\%$;\ a
      first-principles derivation of the substrate $c_J$ is
      deferred (Section~\ref{sec:G3-cutoff-shift}).
\item The Paper~17 identity $\mPl^2/\me^2 = |A_{K_7}|^{-2}$ is
      identified as the structural factor relating the ultraviolet
      cutoff and the infrared electron sector; the $10^{45}$
      Phase~5 gap of Paper~16 is structurally identified with the
      missing $K_7$ Lie-theoretic factor that the bare
      $S^3/2I$ Casimir computation did not include.  A corrected
      Phase~5 calculation incorporating this factor is sketched
      (Section~\ref{sec:G3-phase5-sketch}); the full evaluation
      is logged as a follow-up.
\item Variation of the Sakharov-induced effective action $\Gamma[g]$
      with respect to $g^{\mu\nu}$ yields the full nonlinear
      Einstein equations with cosmological constant and Stelle
      higher-curvature corrections,
      \begin{equation}
        G_{\mu\nu} + \Lambda_{\mathrm{cc}}\,g_{\mu\nu}
        = 8\pi G\,T_{\mu\nu}
        + \alpha_S\,\partial_\mu\partial_\nu R + \cdots,
      \end{equation}
      with $G$ matching Paper~17's structurally derived value.
\item The Schwarzschild metric is verified as a vacuum solution by
      direct symbolic computation: all $10$ independent components
      of the symmetric Ricci tensor $R_{\mu\nu}$ vanish identically,
      evaluated in the \texttt{nwt\_substrate.gravity} module of
      the companion library.
\item Linearisation around flat space reproduces the
      gravitational-wave equation
      $\Box\,h^{\mathrm{TT}}_{\mu\nu} = -16\pi G\,T^{\mathrm{TT}}_{\mu\nu}$
      with two physical transverse-traceless polarizations, and
      the static weak-field limit recovers Newton's law with the
      predicted $G$.
\item Numerically, the resulting prediction agrees with CODATA $G$
      to $+29$\,ppm, inside the $\pm 22$\,ppm experimental error
      bar on $G$ itself.
\end{enumerate}

The structural picture is that gravity in NWT is the macroscopic
projection of a single algebraic object---the $K_7$ graph-state
amplitude---through two complementary mechanisms.  Papers~14--17
projected it as a static structural identity for $G$;
this paper projects it as the dynamical Einstein--Hilbert action
governing arbitrary backgrounds.  The two projections are the same
$K_7$ amplitude read at different scales; the substrate-monism
companion paper~\cite{paper19} assembles this into the broader
master table of NWT observables.

\subsection{Outline}

Section~\ref{sec:matter-lagrangian} sets up the covariant matter
action, identifies its BPS vacuum, and computes the matter
stress-energy tensor $T^{(M)}_{\mu\nu}$ around the vacuum (G1).
Section~\ref{sec:G2-flat-sakharov} performs the matter loop on a
flat background, extracts the Einstein--Hilbert coefficient via the
heat-kernel $a_2$ machinery, and motivates the Sakharov cutoff
choice $\Lambda_{\mathrm{UV}} = \mPl$ (G2).  Section~\ref{sec:G3-K7-bridge},
the technical headline, identifies the $K_7$ amplitude as the
ultraviolet/infrared bridge, derives the
$N_{\mathrm{DOF}} = 12\pi$ structural condition, and closes the
$10^{45}$ Phase~5 gap of Paper~16 (G3).  Section~\ref{sec:G4-linearized}
linearizes the resulting effective action around flat space and
recovers the gravitational-wave equation, the two physical TT
polarizations, and Newton's law in the non-relativistic limit
(G4).  Section~\ref{sec:G5-curved} extends the heat-kernel
expansion to a curved background, picking up Stelle higher-curvature
corrections and a (formally divergent) cosmological-constant
contribution (G5).  Section~\ref{sec:G6-einstein} varies the full
effective action with respect to $g^{\mu\nu}$ and verifies the
resulting nonlinear Einstein equations against Schwarzschild
symbolically (G6).  Section~\ref{sec:open} catalogues open
problems and Section~\ref{sec:conclusion} concludes.

% ==========================================================
\section{The NWT matter Lagrangian at the BPS vacuum}
\label{sec:matter-lagrangian}
% ==========================================================

The Sakharov derivation requires a matter Lagrangian to integrate
over.  In NWT this is the BPS three-field theory of
Paper~16~\cite{paper16}, which we recapitulate here in covariant
form, identify its homogeneous vacuum, and compute the matter
stress-energy tensor that couples to the linearized metric
perturbation.  The full symbolic derivation is implemented in
\texttt{analysis/nwt\_paper18\_G1\_setup.py}; this section
reproduces the analytic content.

\subsection{Field content (L1)}

The minimal NWT matter content is
\begin{itemize}
\item a complex scalar $\psi(x)$, the substrate condensate
      (parametrised by a vacuum expectation value $v$, with
      $|\psi|=v$ in the homogeneous BPS vacuum);
\item a $\U(1)$ gauge field $A_\mu(x)$, gauging the $\psi$ phase
      with charge $e$ and field strength
      $F_{\mu\nu} = \partial_\mu A_\nu - \partial_\nu A_\mu$;
\item an $S^2$-valued Skyrme--Faddeev field $n^a(x)$ ($a=1,2,3$)
      subject to the constraint $n^a n^a = 1$, supporting Hopfion
      windings of knotted-defect topology.
\end{itemize}
This content is derived in Paper~16 from the
$\mathrm{Spin}(7)/\mathrm{Cl}(0,7)/2\mathrm{T}$ substrate algebra:
the complex scalar carries the substrate's $\U(1)_\psi$ phase, the
gauge field gauges that phase (one of the $21$ generators of
$\so(7)$), and the $S^2$-valued $n^a$ realizes the residual
internal $\so(3)$ Cartan of the substrate after the projection to
$3{+}1$ dimensions.  Higher-energy ultraviolet completions
(Paper~14, Paper~16~Phase L1b) extend the gauge sector to
$\so(10)$ at $E_{\mathrm{GUT}} \approx 7.4 \times 10^{15}$\,GeV;
those degrees of freedom enter the Sakharov calculation through
the cutoff (Section~\ref{sec:G3-K7-bridge}) but are not active
infrared modes.

\subsection{Kinetic and gauge structure (L2)}

The covariantly-coupled matter action is
\begin{equation}
\begin{split}
S_{\mathrm{matter}}[\psi, A, n; g] \;=\; \int d^4 x\,\sqrt{-g}\,\Big[\;
    & g^{\mu\nu}(D_\mu\psi)^*(D_\nu\psi) \\
    & - \tfrac{1}{4}\,g^{\mu\alpha} g^{\nu\beta}\,F_{\mu\nu} F_{\alpha\beta} \\
    & - \tfrac{\lambda}{4}\,(|\psi|^2 - v^2)^2 \\
    & + \tfrac{f_\pi^2}{2}\,g^{\mu\nu}\,(\partial_\mu n^a)(\partial_\nu n^a) \\
    & - \tfrac{1}{4 e_{\mathrm{Sk}}^2}\,
       (g^{\mu\alpha} g^{\nu\beta} - g^{\mu\beta} g^{\nu\alpha})\,
       (\partial_\mu n^a)(\partial_\nu n^b)
       (\partial_\alpha n^a)(\partial_\beta n^b)
\;\Big],
\end{split}
\label{eq:Smatter-covariant}
\end{equation}
with $D_\mu \psi = (\partial_\mu - i e A_\mu)\psi$.  The four
free parameters are $\{e, v, \lambda, f_\pi, e_{\mathrm{Sk}}\}$, of
which Paper~16 fixes $\lambda$ at the \emph{BPS critical coupling}
\begin{equation}
  \lambda \;=\; \frac{e^2}{2},
  \label{eq:BPS-coupling}
\end{equation}
which equates the Higgs and gauge inverse-correlation lengths so
that vortex Bogomolny bounds are saturated at first order.  The
remaining parameters are gauged-sector parameters whose values
enter Paper~6's mass spectrum but, importantly for our purposes
here, drop out of the Sakharov $G$-coefficient at flat-space leading
order (Section~\ref{sec:G2-flat-sakharov}).

\subsection{The topological term (L3)}

The full L3 sector adds the topological Hopf charge
\begin{equation}
  S_{\mathrm{top}}\;=\;\theta\int d^4 x\,Q_H[n],
\end{equation}
where $Q_H$ is the integer Hopf invariant of $n: \R^3 \to S^2$.
Two facts make this term invisible to the Sakharov derivation:
$Q_H$ depends only on the topological winding of $n$, not on the
metric, so $\delta S_{\mathrm{top}}/\delta g^{\mu\nu} = 0$; and the
homogeneous BPS vacuum has $Q_H = 0$, so the term carries no
energy at the vacuum point.  $S_{\mathrm{top}}$ does contribute
crucially to the matter spectrum (Paper~6's mass formula uses
$Q_H = p\cdot m$ for the $(p,q)$ torus-knot windings), but it
generates no graviton couplings.  We therefore drop it from the
Sakharov calculation that follows.

\subsection{Linearisation and the matter stress-energy tensor (G1)}
\label{sec:G1}

We linearize around flat space,
\begin{equation}
  g_{\mu\nu} \;=\; \eta_{\mu\nu} + h_{\mu\nu},
  \qquad
  g^{\mu\nu} \;=\; \eta^{\mu\nu} - h^{\mu\nu} + \mathcal{O}(h^2),
  \qquad
  \sqrt{-g} \;=\; 1 + \tfrac{1}{2}\,h + \mathcal{O}(h^2),
  \label{eq:linearization}
\end{equation}
with $h \equiv \eta^{\mu\nu} h_{\mu\nu}$.  Substituting
\eqref{eq:linearization} into \eqref{eq:Smatter-covariant} and
collecting first-order terms in $h$ yields the standard form
\begin{equation}
  S_{\mathrm{matter}}[\psi, A, n; \eta+h]
  \;=\;
  S_{\mathrm{matter}}^{(0)}[\psi, A, n]
  \;-\;
  \tfrac{1}{2}\int d^4 x\,h^{\mu\nu}\,T^{(M)}_{\mu\nu}[\psi, A, n]
  \;+\;\mathcal{O}(h^2),
  \label{eq:Smatter-linear}
\end{equation}
where $T^{(M)}_{\mu\nu}$ is the matter stress-energy tensor in
flat space.  The Sakharov mechanism is then the matter loop on
\eqref{eq:Smatter-linear}: integrating out the matter fields
generates an effective action for $h$.

\paragraph{The homogeneous BPS vacuum.}  The minimum of
$S_{\mathrm{matter}}^{(0)}$ at constant fields is
\begin{equation}
  \psi_0 = v,
  \qquad
  A_\mu = 0,
  \qquad
  n^a_0 = \delta^{a 3}.
  \label{eq:BPS-vacuum}
\end{equation}
At this point the Higgs potential vanishes,
$|\psi_0|^2 = v^2 \Rightarrow (|\psi_0|^2 - v^2)^2 = 0$; all
gradients are zero so the kinetic terms vanish; and the Skyrme
quartic term is fourth order in gradients, also zero.  The
flat-space Lagrangian density and the matter stress-energy tensor
both vanish at the vacuum:
\begin{equation}
  \mathcal{L}_{\mathrm{flat}}^{\mathrm{vac}} = 0,
  \qquad
  T^{(M)}_{\mu\nu}\big|_{\mathrm{vac}} = 0.
  \label{eq:vac-vanish}
\end{equation}
The matter sector therefore contributes \emph{no classical
cosmological constant} at the BPS vacuum.  Quantum loops at the
cutoff scale $\Lambda_{\mathrm{UV}}$ will generate one of the form
$\Lambda_{\mathrm{cc}} \sim \Lambda_{\mathrm{UV}}^4 / (16\pi G)$
(Sakharov's original observation); this is the standard
cosmological-constant problem and we track it separately in
Section~\ref{sec:G5-curved}, deferring its detailed treatment to
Paper~21.

\paragraph{Second-order fluctuations.}  Around \eqref{eq:BPS-vacuum},
we parametrise
\begin{equation}
  \psi = v + \phi_1 + i\phi_2,
  \qquad
  n^a = \big(n_1,\,n_2,\,1 - \tfrac{1}{2}(n_1^2 + n_2^2) + \mathcal{O}(n^4)\big),
  \label{eq:fluct-param}
\end{equation}
where $(\phi_1, \phi_2)$ are the radial and Goldstone fluctuations
of $\psi$, and $(n_1, n_2)$ are the two $S^2$ Goldstones around the
north pole.  At second order in fluctuations, the flat-space
Lagrangian splits into kinetic and mass terms for each fluctuation
type:
\begin{equation}
\mathcal{L}_{\mathrm{flat}}^{(2)}
  = (\partial_\mu \phi_1)^2
  + (\partial_\mu \phi_2 - e v A_\mu)^2
  - \lambda v^2\,\phi_1^2
  + \tfrac{f_\pi^2}{2}\big[(\partial_\mu n_1)^2 + (\partial_\mu n_2)^2\big]
  - \tfrac{1}{4}\,F_{\mu\nu}F^{\mu\nu}.
\label{eq:Lflat-quadratic}
\end{equation}
The $(\phi_2 - e v A_\mu)$ combination is the standard Higgs
mechanism eating: $\phi_2$ is the Goldstone of the broken $\U(1)$
and is absorbed into a longitudinal polarization of $A_\mu$.  The
radial mode $\phi_1$ acquires a mass $m_\psi^2 = \lambda v^2$.  The
two $S^2$ Goldstones $n_1, n_2$ are massless at this order
(Skyrme kinetic only); their masses arise from the Hopf term and
substrate corrections that we have dropped here.

\paragraph{The matter stress-energy.}  Direct variation of
\eqref{eq:Lflat-quadratic} with respect to $g^{\mu\nu}$ at
$g = \eta$ (or, equivalently, evaluation of the canonical
energy-momentum tensor) yields, sector by sector,
\begin{align}
T^{\psi}_{\mu\nu}
  &= 2\,(\partial_\mu \phi_1)(\partial_\nu \phi_1)
  + 2\,(\partial_\mu \phi_2 - e v A_\mu)(\partial_\nu \phi_2 - e v A_\nu)
  - \eta_{\mu\nu}\,\mathcal{L}^{(2)}_\psi,
  \label{eq:Tpsi}\\
T^{A}_{\mu\nu}
  &= F_{\mu\alpha}\,{F_\nu}^\alpha
  - \tfrac{1}{4}\,\eta_{\mu\nu}\,F_{\alpha\beta} F^{\alpha\beta},
  \label{eq:TA}\\
T^{n}_{\mu\nu}
  &= f_\pi^2\big[(\partial_\mu n_1)(\partial_\nu n_1)
                + (\partial_\mu n_2)(\partial_\nu n_2)\big]
  - \eta_{\mu\nu}\,\mathcal{L}^{(2)}_n.
  \label{eq:Tn}
\end{align}
The total is $T^{(M)}_{\mu\nu} = T^{\psi}_{\mu\nu} + T^{A}_{\mu\nu}
+ T^{n}_{\mu\nu}$.  Each piece is bilinear in matter fluctuations
(the BPS vacuum cancels at the linear order, by
\eqref{eq:vac-vanish}), so the matter--graviton coupling
$-\tfrac{1}{2} h^{\mu\nu} T^{(M)}_{\mu\nu}$ produces a
graviton-matter-matter vertex with two matter legs: this is
\emph{the vertex Sakharov's mechanism requires} for the matter loop
in Section~\ref{sec:G2-flat-sakharov}.

\paragraph{Trace identity.}  As a sanity check, the trace
$T^{(M){\,\mu}}{}_\mu = \eta^{\mu\nu} T^{(M)}_{\mu\nu}$ collects
into
\begin{equation}
  T^{(M){\,\mu}}{}_\mu
  \;=\; -2\,\mathcal{L}^{(2)}_\psi - 2\,\mathcal{L}^{(2)}_n
        + (D_\psi - 4)\,\big[\text{kinetic terms}\big]
        + 4\,\big[\text{mass terms}\big],
\end{equation}
which in $D=4$ vanishes for the kinetic-only contributions and
leaves the explicit mass term $-2 \lambda v^2 \phi_1^2$
(corresponding to the classical conformal-anomaly contribution from
the radial $\psi$ mode).  This matches the textbook expectation
for a massive scalar in four dimensions and confirms the absence
of trace anomalies at the classical level.

The output of this section---the matter--graviton vertex
$-\tfrac{1}{2} h^{\mu\nu} T^{(M)}_{\mu\nu}$ with $T^{(M)}_{\mu\nu}$
given by \eqref{eq:Tpsi}--\eqref{eq:Tn}---is the integrand for the
matter loop performed in Section~\ref{sec:G2-flat-sakharov}.

% ==========================================================
\section{Sakharov-induced gravity on a flat background (G2)}
\label{sec:G2-flat-sakharov}
% ==========================================================

With the matter--graviton vertex \eqref{eq:Smatter-linear} in hand,
the Sakharov program proceeds by integrating out the matter
fluctuations $\{\phi_1, \phi_2, A_\mu, n_1, n_2\}$ at one loop on
$g_{\mu\nu} = \eta_{\mu\nu} + h_{\mu\nu}$ and reading off the
coefficient of $\sqrt{-g}\,R$ in the resulting effective action for
$h$.  The full numerical evaluation is in
\texttt{analysis/nwt\_paper18\_G2\_matter\_loop.py}.

\subsection{The matter loop and the heat-kernel expansion}

For each matter species with second-order kinetic operator
$K = -\Box + m^2 + \xi R$ on a curved background, the standard
one-loop effective action is~\cite{vassilevich2003}
\begin{equation}
  \Gamma_{\text{1-loop}}[g]
   \;=\;
   -\,\tfrac{1}{2}\,\Tr \ln K
   \;=\;
   -\,\tfrac{1}{2}\,\int_0^\infty \frac{dt}{t}\,\Tr\!\bigl[\,e^{-t K}\bigr],
  \label{eq:gamma-trln}
\end{equation}
where the proper-time integral is regulated at the lower limit by
$t_{\min} \sim 1/\Lambda_{\mathrm{UV}}^2$.  Inserting the
Seeley--DeWitt heat-kernel expansion
\begin{equation}
  \Tr\!\bigl[\,e^{-t K}\bigr]
   \;=\;
   \frac{1}{(4\pi t)^{d/2}}\,
   \int d^d x\,\sqrt{g}\;\bigl[\,a_0(x) + a_1(x)\,t + a_2(x)\,t^2 + \cdots\,\bigr],
  \label{eq:hk-expansion}
\end{equation}
and integrating in $d = 4$ with cutoff $\Lambda_{\mathrm{UV}}$,
\eqref{eq:gamma-trln} becomes
\begin{equation}
  \Gamma_{\text{1-loop}}[g]
   \;=\;
   -\,\frac{1}{32\pi^2}\,\int d^4 x\,\sqrt{g}\,
   \Big[\,
       \tfrac{1}{2}\,\Lambda_{\mathrm{UV}}^4\,a_0
     + \Lambda_{\mathrm{UV}}^2\,a_1
     + \ln(\Lambda_{\mathrm{UV}}^2/m^2)\,a_2
     + \mathrm{finite}
   \,\Big].
  \label{eq:gamma-cutoff}
\end{equation}
For a minimally coupled massive scalar ($\xi = 0$),
\begin{align}
  a_0 &= 1, \\
  a_1 &= \tfrac{1}{6}\,R - m^2, \\
  a_2 &= \tfrac{1}{180}\bigl(R^2 - R_{\mu\nu} R^{\mu\nu} + R_{\mu\nu\rho\sigma} R^{\mu\nu\rho\sigma}\bigr)
       + \tfrac{1}{2}\bigl(m^2 - \tfrac{R}{6}\bigr)^2
       + \tfrac{1}{30}\,\Box R + \cdots,
\end{align}
all evaluated at the field point $x$.

\subsection{Extracting the Einstein--Hilbert coefficient}

The induced Einstein--Hilbert action is the long-wavelength,
single-power-of-curvature piece of \eqref{eq:gamma-cutoff}:
\begin{equation}
  \Gamma_{\mathrm{EH}}[g]
   \;=\;
   -\,\int d^4 x\,\sqrt{g}\;
       \frac{1}{16\pi G_{\mathrm{ind}}}\,R(g),
  \label{eq:gamma-EH}
\end{equation}
which by direct comparison to the $a_1$ contribution of
\eqref{eq:gamma-cutoff} (the $R/6$ term in the
$\Lambda_{\mathrm{UV}}^2$ coefficient) gives, per scalar species,
\begin{equation}
  \frac{1}{16\pi G_{\mathrm{ind}}}\bigg|_{\text{1 scalar}}
   \;=\;
   \frac{\Lambda_{\mathrm{UV}}^2}{192\pi^2}
   \;-\;\frac{m^2}{192\pi^2}\,\ln\!\bigl(\Lambda_{\mathrm{UV}}^2/m^2\bigr)
   \;+\;\cdots
  \label{eq:per-scalar}
\end{equation}
The full matter-sector Sakharov coefficient is the sum of
\eqref{eq:per-scalar} over the matter species, weighted by their
spin:
\begin{equation}
  \frac{1}{16\pi G_{\mathrm{ind}}}
   \;=\;
   \sum_{\text{species}}\,c_J^{\text{species}}\;
   \frac{\Lambda_{\mathrm{UV}}^2}{192\pi^2}
   \;+\;\mathcal{O}(m^2 \ln \Lambda).
  \label{eq:total-sakharov}
\end{equation}
The standard spin coefficients (Birrell--Davies, Vassilevich) are
$c_0 = +1$ for each real scalar and $c_1 = +3$ for each massive
Higgsed vector (Faddeev--Popov ghosts already absorbed).

\paragraph{NWT matter-DOF count.}  At the BPS vacuum
\eqref{eq:BPS-vacuum}, the fluctuation spectrum of L1+L2 consists of
\begin{table}[H]
\centering
\begin{tabular}{lcrl}
\toprule
\textbf{species} & \textbf{spin} & \textbf{DOF} & \textbf{mass} \\
\midrule
$\phi_1$ (radial Higgs)                        & $0$ & $1$ & $m_\psi^2 = \lambda v^2$ \\
$A_\mu$ (Higgsed vector, $\phi_2$ eaten)       & $1$ & $3$ & $m_A^2 = e^2 v^2$ \\
$n_1, n_2$ ($S^2$ Goldstones)                  & $0$ & $2$ & $0$ \\
\midrule
$N_{\mathrm{DOF}} = 1 + 3 + 2$                 &     & $\mathbf{6}$ & \\
\bottomrule
\end{tabular}
\caption{Bosonic degrees of freedom of the BPS NWT matter sector
contributing to the Sakharov $\Lambda^2$ term.  The L3 Hopf charge
contributes no propagating modes; the Skyrme quartic is a derivative
self-interaction that enters only at higher loop order.}
\label{tab:bps-dof}
\end{table}

The leading-order Sakharov $1/G$ coefficient is therefore
\begin{equation}
  \boxed{\;
    \frac{1}{16\pi G_{\mathrm{ind}}}^{\!\!\text{LO,\;BPS}}
     \;=\;
     \frac{N_{\mathrm{DOF}}\,\Lambda_{\mathrm{UV}}^2}{192\pi^2}
     \;=\;
     \frac{6\,\Lambda_{\mathrm{UV}}^2}{192\pi^2}
     \;=\;
     \frac{\Lambda_{\mathrm{UV}}^2}{32\pi^2}.
  \;}
  \label{eq:sakharov-NWT}
\end{equation}
This is the Sakharov $G$ from the $6$-DOF BPS sector alone,
expressed in terms of an as-yet-unspecified cutoff
$\Lambda_{\mathrm{UV}}$.

\subsection{The Sakharov cutoff and the choice
$\Lambda_{\mathrm{UV}} = \mPl$}

In the original Sakharov picture, $\Lambda_{\mathrm{UV}}$ is a
free parameter to be fixed by some external input.  In NWT it is
\emph{determined}: equating \eqref{eq:sakharov-NWT} with Paper~17's
structural prediction \eqref{eq:p17}, and using
$1/(16\pi G) = \mPl^2/(16\pi)$ in natural units, gives
\begin{equation}
  \frac{N_{\mathrm{DOF}}\,\Lambda_{\mathrm{UV}}^2}{192\pi^2}
   \;=\;
   \frac{\mPl^2}{16\pi}
   \;=\;
   \frac{1}{16\pi}\,
       \frac{\me^2}{(8/7)^2 (1 + \alpha/7)^2\,\alpha^{21}}.
  \label{eq:cutoff-equation}
\end{equation}
Solving for $\Lambda_{\mathrm{UV}}$:
\begin{equation}
  \Lambda_{\mathrm{UV}}^2
   \;=\;
   \frac{12\pi}{N_{\mathrm{DOF}}}\,\mPl^2,
  \qquad
  \Lambda_{\mathrm{UV}}
   \;=\;
   \mPl\,\sqrt{\frac{12\pi}{N_{\mathrm{DOF}}}}.
  \label{eq:cutoff-solved}
\end{equation}
For the BPS-sector $N_{\mathrm{DOF}} = 6$, this yields
$\Lambda_{\mathrm{UV}} \approx \mPl\,\sqrt{2\pi}$, a factor
$\sqrt{2\pi} \approx 2.51$ above the Planck scale.  This factor is
the gap Paper~16 reported as Phase~5 and deferred:\ Sakharov's naive
$\Lambda^2$ scaling alone does not reproduce Paper~17's
structurally-derived $G$ if $\Lambda_{\mathrm{UV}} = \mPl$
identically.

\paragraph{Inverting the question.}  The clean way to read
\eqref{eq:cutoff-solved} is structural rather than numerical.
\emph{If} $\Lambda_{\mathrm{UV}} = \mPl$ (as standard
Sakharov-induced gravity treatments take it), then the matter-DOF
count required for the Sakharov coefficient to land on Paper~17's
$1/G$ is
\begin{equation}
  \boxed{\;\;
    N_{\mathrm{DOF}}^{\mathrm{required}}
     \;=\;
     12\pi
     \;\approx\;
     37.70.
  \;\;}
  \label{eq:N-required}
\end{equation}
The minimal BPS sector contributes $6$; the GUT extension of
Paper~14 supplies $33$ additional bosonic degrees of freedom in
the $\so(10)$ adjoint at $E_{\mathrm{GUT}} \approx 7.4 \times 10^{15}$
GeV~\cite{paper14}.  The total
\begin{equation}
  N_{\mathrm{DOF}}^{\mathrm{GUT}}
   \;=\;
   N_{\mathrm{DOF}}^{\mathrm{BPS}} + N_{\mathrm{DOF}}^{\so(10)}
   \;=\;
   6 + 33
   \;=\;
   39
  \label{eq:N-GUT}
\end{equation}
matches the required $12\pi$ to $3.4\%$, consistent with the
loop-uncomputed contributions from interactions and the
$\mathcal{O}(m^2 \ln \Lambda)$ corrections in
\eqref{eq:per-scalar}.  We will return to the precise structural
status of this matching in
Section~\ref{sec:G3-K7-bridge}, where the $K_7$ amplitude is
identified as the bridge that closes the residual gap.

The takeaway from this section is that the leading-order
Sakharov-induced $G$ from the BPS matter sector alone is
\emph{too small} to reproduce CODATA $G$ by a factor
$\sqrt{2\pi}$ in the cutoff (or, equivalently, the matter-DOF count
is too small by a factor $2\pi$).  Closing this gap requires
either (i) extending the matter content to the GUT
completion---which gets within $3.4\%$---or (ii) the additional
structural factor that the $K_7$ Wilson amplitude provides.  The
correct view, established next, is that these are not alternatives
but the same fact at two levels of resolution.

% ==========================================================
\section{The $K_7$ amplitude bridges UV and IR (G3) -- HEADLINE}
\label{sec:G3-K7-bridge}
% ==========================================================

This is the technical headline of the paper.  Section~\ref{sec:G2-flat-sakharov}
showed that the bare Sakharov calculation with the BPS matter
sector requires $N_{\mathrm{DOF}} = 12\pi$ for $\Lambda_{\mathrm{UV}}
= \mPl$.  Here we show that this result is not a free-parameter
adjustment but the precise statement that the $K_7$ Wilson amplitude
of Paper~17 is the structural factor bridging the matter-loop
ultraviolet and the electron-sector infrared.  The full numerical
verification is in
\texttt{analysis/nwt\_paper18\_G3\_K7\_sakharov.py}.

\subsection{Two independent routes to $1/G$}
\label{sec:G3-two-routes}

NWT now produces $1/G$ along two independent paths:

\paragraph{Route A (Sakharov, this paper).}  From the BPS
matter loop on a flat background \eqref{eq:sakharov-NWT},
\begin{equation}
  \left[\frac{1}{16\pi G}\right]_{\mathrm{Sakharov}}
   \;=\;
   \frac{N_{\mathrm{DOF}}\,\Lambda_{\mathrm{UV}}^2}{192\pi^2}.
  \label{eq:routeA}
\end{equation}
Inputs:\ the matter-DOF count $N_{\mathrm{DOF}}$ and the cutoff
$\Lambda_{\mathrm{UV}}$.

\paragraph{Route B ($K_7$ amplitude, Paper~17).}  From the
$\mathrm{Spin}(7)$ Chern--Simons Wilson amplitude on the $K_7$ graph
on the Heegaard torus of $S^3/2I$~\cite{paper17},
\begin{equation}
  \frac{\me}{\mPl}
   \;=\;
   \mathrm{bracket}(\alpha)\;\alpha^{21/2},
  \qquad
  \mathrm{bracket}(\alpha)
   \;\equiv\;
   \frac{\dim S}{\dim V}\,
   \left[\,1 + \frac{\alpha}{\dim V} + \frac{\dim \Adj}{\dim S}\,\alpha^2\,\right]
   \;=\;
   \tfrac{8}{7}\bigl(1 + \tfrac{\alpha}{7} + 3\alpha^2\bigr),
  \label{eq:bracket-def}
\end{equation}
which (in natural units $\hbar = c = 1$) gives
$G\,\me^2 = \mathrm{bracket}^2\,\alpha^{21}$ and
\begin{equation}
  \left[\frac{1}{16\pi G}\right]_{K_7}
   \;=\;
   \frac{1}{16\pi\,\mathrm{bracket}(\alpha)^2\,\alpha^{21}\,/\,\me^2}
   \;=\;
   \frac{\me^2}{16\pi\,\mathrm{bracket}(\alpha)^2\,\alpha^{21}}.
  \label{eq:routeB}
\end{equation}
Inputs:\ the fine-structure constant $\alpha$ and the electron
mass $\me$.  The dimensions $\dim V = 7$, $\dim S = 8$,
$\dim\Adj = 21$ are the $\mathrm{Spin}(7)$ representation theory
that gives \eqref{eq:bracket-def} structural origin (Paper~17 \S6).

\paragraph{The two routes are different physical statements.}
Route A is dynamical: it asks what coefficient appears in front of
$\sqrt{-g}\,R$ when matter is integrated out.  Route B is
structural: it asks what value of $\me/\mPl$ is consistent with the
$\mathrm{Spin}(7)$ Wilson amplitude on $K_7$.  The two routes share
no derivation; consistency between them is a non-trivial constraint
on the matter spectrum and the cutoff.  In what follows we equate
them as a consistency condition (not as two independent
computations of the same quantity) and read off the cutoff in terms
of Paper~17's IR observables, with the matter-content question
that this leaves open addressed in Section~\ref{sec:G3-cutoff-shift}.

\subsection{Equating the two routes}

Setting \eqref{eq:routeA} equal to \eqref{eq:routeB},
\begin{equation}
  \frac{N_{\mathrm{DOF}}\,\Lambda_{\mathrm{UV}}^2}{192\pi^2}
   \;=\;
   \frac{\me^2}{16\pi\,\mathrm{bracket}^2\,\alpha^{21}},
  \label{eq:match-eq}
\end{equation}
and solving for $\Lambda_{\mathrm{UV}}$,
\begin{equation}
  \Lambda_{\mathrm{UV}}
   \;=\;
   \sqrt{\frac{12\pi}{N_{\mathrm{DOF}}}}\;
   \frac{\me}{\mathrm{bracket}(\alpha)\,\alpha^{21/2}}
   \;=\;
   \sqrt{\frac{12\pi}{N_{\mathrm{DOF}}}}\;\mPl,
  \label{eq:lambda-uv-K7}
\end{equation}
where the second equality uses Paper~17's definition
$\mPl = \me / [\mathrm{bracket}(\alpha)\,\alpha^{21/2}]$.

The structurally clean statement is therefore:
\begin{equation}
  \boxed{\;
    \Lambda_{\mathrm{UV}}
     \;=\;
     \mPl
    \quad\Longleftrightarrow\quad
    N_{\mathrm{DOF}}
     \;=\;
     12\pi\;\approx\;37.70.
  \;}
  \label{eq:central-identity}
\end{equation}

Equation \eqref{eq:central-identity} is a consistency relation:
exact equality $\Lambda_{\mathrm{UV}} = \mPl$ holds iff
$N_{\mathrm{DOF}}$ takes the transcendental value $12\pi$, which no
integer-valued substrate matter content can deliver.  The
generalized relation $\Lambda_{\mathrm{UV}}/\mPl =
\sqrt{12\pi/N_{\mathrm{DOF}}}$ is the operationally useful form;\
\eqref{eq:central-identity} expresses the limiting case in compact
notation.  Section~\ref{sec:G3-cutoff-shift} addresses what
matter content can saturate this relation to within a few percent.

\subsection{The structural reading: $\Lambda_{\mathrm{UV}}$ is shifted
            from $\mPl$ by a calculable group-theoretic factor}
\label{sec:G3-cutoff-shift}

The matching condition \eqref{eq:lambda-uv-K7} is most cleanly read as
a relation between an integer-valued matter input and a single
dimensionless cutoff ratio:
\begin{equation}
  \boxed{\;
    \frac{\Lambda_{\mathrm{UV}}}{\mPl}
     \;=\;
     \sqrt{\frac{12\pi}{N_{\mathrm{DOF}}}}.
  \;}
  \label{eq:cutoff-shift}
\end{equation}
The $12\pi = 192\pi^2/(16\pi)$ prefactor is fixed by the heat-kernel
normalization of \eqref{eq:routeA}; $N_{\mathrm{DOF}}$ is the
$c_J$-weighted matter content of \eqref{eq:total-sakharov}.  Two
natural choices arise within the NWT hierarchy:
\begin{enumerate}
\item The minimal BPS sector of Section~\ref{sec:G2-flat-sakharov},
      $N_{\mathrm{DOF}}^{\mathrm{BPS}} = 6$, giving
      $\Lambda_{\mathrm{UV}}/\mPl = \sqrt{2\pi} \approx 2.51$.
      This places the cutoff above $\mPl$, which is unphysical in the
      Sakharov framework: the BPS sector alone is therefore
      incomplete.
\item The $\so(10)$ GUT extension of Paper~14, in which the matter
      content is supplemented by the gauge bosons of the
      $\so(10)/\mathrm{SM}$ coset.  Counting these requires care; the
      result depends on whether one counts coset \emph{generators} or
      \emph{propagating bosonic degrees of freedom}, which differ by
      the spin multiplicity of the broken gauge bosons.  We address
      this counting next.
\end{enumerate}

\paragraph{Counting the GUT contribution to $N_{\mathrm{DOF}}$.}
$\so(10)$ has $\dim\so(10) = 45$ generators, of which the Standard
Model gauge group $\SU(3) \times \SU(2) \times \U(1)$ accounts for
$8 + 3 + 1 = 12$.  The coset
$\so(10)/(\SU(3) \times \SU(2) \times \U(1))$ has
\begin{equation}
  \dim_{\mathrm{coset}}
   \;=\;
   45 - 12
   \;=\;
   33
\end{equation}
generators.  In the Higgsed phase below $E_{\mathrm{GUT}}$, each
coset generator is realized as a massive vector boson with three
propagating polarizations (transverse plus longitudinal-from-eaten
Goldstone), so the propagating-DOF count is
\begin{equation}
  N_{\mathrm{DOF}}^{\so(10)\,\mathrm{prop}}
   \;=\;
   3 \times 33
   \;=\;
   99,
  \label{eq:N-prop-naive}
\end{equation}
not $33$.  Combined with the BPS sector,
$N_{\mathrm{DOF}}^{\mathrm{BPS}} + N_{\mathrm{DOF}}^{\so(10)\,\mathrm{prop}}
= 6 + 99 = 105$, which fails to match $12\pi \approx 37.70$ by a
factor of nearly three.  The naive Sakharov sum over all
SM-and-GUT propagating DOFs is therefore \emph{not} the correct
input to \eqref{eq:cutoff-shift}.

\paragraph{The decoupling resolution.}
The standard resolution is that the heavy GUT bosons decouple from
the low-energy effective theory at a Sakharov cutoff
$\Lambda_{\mathrm{UV}} \sim \mPl \gg E_{\mathrm{GUT}}$.  Their
contribution to the matter-loop $a_1$ coefficient comes with the
Appelquist--Carazzone~\cite{appelquist1975} suppression factor
$\sim m_W^{\mathrm{GUT}\,2}/\Lambda_{\mathrm{UV}}^2$, where
$m_W^{\mathrm{GUT}} \sim E_{\mathrm{GUT}} \approx
7.4 \times 10^{15}\,\GeV$ is the heavy-vector mass scale and
$\Lambda_{\mathrm{UV}} \approx \mPl \approx 1.2 \times 10^{19}\,\GeV$.
The suppression is
\begin{equation}
  \frac{m_W^{\mathrm{GUT}\,2}}{\mPl^2}
   \;\approx\;
   \left(\frac{7.4 \times 10^{15}}{1.2 \times 10^{19}}\right)^{\!2}
   \;\approx\;
   3.8 \times 10^{-7},
\end{equation}
i.e.\ each heavy GUT vector contributes effectively $\sim 10^{-6}$
to $N_{\mathrm{DOF}}$ at IR scales, not $3$.  The GUT bosons are
present in the spectrum but \emph{absent from the IR Sakharov
coefficient at leading order}.

\paragraph{What does count: substrate fields, not IR projections.}
The Sakharov calculation of Section~\ref{sec:G2-flat-sakharov}
integrates out the matter Lagrangian of Paper~16:\ the complex
scalar $\psi$, the $\U(1)$ gauge field $A_\mu$ at the BPS coupling,
and the Skyrme--Faddeev field $n^a$ on $S^2$.  This is the
\emph{substrate} matter content.  It contains no fundamental Dirac
fermions and no fundamental massless gauge bosons; the IR-projected
SM fermions, photon, and gluons that we observe at low energy
appear as topological projections of the substrate (vortex
configurations carrying $(p,q)$-knot quantum numbers, K\textsubscript{7}
graph-state content, $\mathrm{Spin}(7)$ Wilson amplitudes), not as
new fields entering the heat kernel of \eqref{eq:total-sakharov}.

This distinction is the single most important point of this
subsection.  Identifying $N_{\mathrm{DOF}}$ in
\eqref{eq:cutoff-shift} requires summing $c_J$ over the
\emph{substrate} fields only;\ adding a separate IR-projected
contribution would double-count the matter and conflate two
different levels of description.

\paragraph{What goes wrong if you forget the distinction.}
A naive Birrell--Davies sum over the IR-projected SM spectrum
(without recognizing it as projected substrate content) would give:
\begin{table}[H]
\centering
\begin{tabular}{lrrr}
\toprule
\textbf{IR species (illustrative; not what NWT sums)} & \textbf{count} & \textbf{$c_J$} & \textbf{contribution} \\
\midrule
$W^\pm, Z$ massive vectors                       & $3$ & $+3$ & $+9$ \\
photon $\gamma$ massless vector                  & $1$ & $-2$ & $-2$ \\
gluons $g$ massless                              & $8$ & $-2$ & $-16$ \\
SM Higgs scalar                                  & $1$ & $+1$ & $+1$ \\
SM Dirac fermions ($\sim 22.5$ Dirac, 3 gen.)    & --  & $-2$ & $\sim -45$ \\
\midrule
\midrule
$N_{\mathrm{DOF}}^{\mathrm{IR,naive}}$           &     &     & $\sim -53$ \\
\bottomrule
\end{tabular}
\caption{Counterexample sum.  Standard Birrell--Davies
coefficients applied to the IR-projected SM spectrum yield a
sharply negative $N_{\mathrm{DOF}}$, which would make
\eqref{eq:cutoff-shift} ill-defined.  This is the well-known
Sakharov fermion-sign issue~\cite{adler1982, visser2002} of
treating the IR spectrum as if its species were fundamental
substrate fields.  In NWT it is a category error:\ IR fermions and
massless gauge bosons are projected substrate configurations, not
new fields.}
\label{tab:N-IR-counterexample}
\end{table}

\paragraph{The substrate-level count.}
The actual $N_{\mathrm{DOF}}$ entering \eqref{eq:cutoff-shift} is
the $c_J$-weighted sum over substrate fields only:
\begin{table}[H]
\centering
\begin{tabular}{lrrr}
\toprule
\textbf{substrate sector} & \textbf{counting} & \textbf{$c_J^{\mathrm{sub}}$} & \textbf{contribution} \\
\midrule
BPS L1+L2 (\S\ref{sec:G2-flat-sakharov})         &     &     & $+6$ \\
$\quad\phi_1$ radial Higgs (scalar)              & $1$ & $+1$ & \\
$\quad A_\mu^{\U(1)}$ Higgsed vector             & $1$ & $+3$ & \\
$\quad n_1, n_2$ Goldstones (scalar)             & $2$ & $+1$ & \\
\midrule
$\so(10)/\mathrm{SM}$ coset (Paper~14, Phase L1b) &    &     & $+33$ or $+99$ \\
$\quad$ generator counting (1 substrate mode/gen) & $33$ & $+1$ & ($+33$) \\
$\quad$ polarization counting (Proca-like)        & $33$ & $+3$ & ($+99$) \\
\midrule
\midrule
$N_{\mathrm{DOF}}^{\mathrm{substrate}}$           &     &     & $39$ or $105$ \\
\bottomrule
\end{tabular}
\caption{Substrate-level $N_{\mathrm{DOF}}$.  All contributions are
positive (the substrate is bosonic at this order).  The ambiguity
between generator counting and polarization counting reflects an
open question about the substrate $c_J$ for the heavy
$\so(10)/\mathrm{SM}$ coset modes:\ each coset generator contributes
$1$ substrate mode (the Cartan-graded element of the Wilson loop)
or $3$ propagating polarizations (Proca-like), depending on whether
the substrate $c_J^{\mathrm{sub}}$ is identified with a single
algebraic mode or with a polarization count.  Heavy GUT modes are
decoupled by Appelquist--Carazzone at IR scales but contribute
unsuppressed at $\Lambda_{\mathrm{UV}} \sim \mPl$.}
\label{tab:N-substrate}
\end{table}

\paragraph{Reading the substrate count against \eqref{eq:cutoff-shift}.}
Inserting the two substrate-level counts into \eqref{eq:cutoff-shift}
gives:
\begin{align}
  \text{generator counting:} & \quad
    N_{\mathrm{DOF}}^{\mathrm{sub}} = 39
    \;\Longrightarrow\;
    \frac{\Lambda_{\mathrm{UV}}}{\mPl}
    = \sqrt{\tfrac{12\pi}{39}}
    \approx 0.983,
    \label{eq:cutoff-39}\\
  \text{polarization counting:} & \quad
    N_{\mathrm{DOF}}^{\mathrm{sub}} = 105
    \;\Longrightarrow\;
    \frac{\Lambda_{\mathrm{UV}}}{\mPl}
    = \sqrt{\tfrac{12\pi}{105}}
    \approx 0.599.
    \label{eq:cutoff-105}
\end{align}
The generator-counting result \eqref{eq:cutoff-39} places
$\Lambda_{\mathrm{UV}}$ within $1.7\%$ of $\mPl$, consistent with
the heuristic ``$N_{\mathrm{DOF}} = 39$ matches $12\pi$ to $3.4\%$''
coincidence reported in earlier drafts.  The polarization-counting
result \eqref{eq:cutoff-105} places it at $\sim 0.6\,\mPl$.  Both
are bosonic, both are positive-signed, both are admissible \emph{a
priori}; selecting between them requires a first-principles
derivation of $c_J^{\mathrm{sub}}$ for the substrate-projected
heavy coset bosons.

\paragraph{The open question: substrate $c_J$.}
The substrate-projection mechanism of the substrate-monism
framework (Paper~19~\cite{paper19}) maps Cartan-graded
algebra modes of the substrate onto IR Proca/Maxwell fields by a
specific (currently underived) projection step.  The corresponding
$c_J^{\mathrm{sub}}$ is the heat-kernel coefficient of the
substrate operator, not the IR Proca operator.  We have not
computed it.  Within the present paper we record the empirical
observation that generator counting (one substrate mode per coset
generator, $c_J^{\mathrm{sub}} = +1$) reproduces the
$N_{\mathrm{DOF}} \approx 12\pi$ matching to $3.4\%$, while
polarization counting ($c_J^{\mathrm{sub}} = +3$) does not.  A
substrate-level heat-kernel computation determining
$c_J^{\mathrm{sub}}$ from first principles is logged as a follow-up
technical note.

\paragraph{The structurally clean statement.}
The Sakharov cutoff in NWT is the geometric mean
$\Lambda_{\mathrm{UV}} = \mPl\,\sqrt{12\pi/N_{\mathrm{DOF}}^{\mathrm{sub}}}$,
with the $12\pi$ factor fixed by the heat-kernel normalization and
$N_{\mathrm{DOF}}^{\mathrm{sub}}$ a $c_J^{\mathrm{sub}}$-weighted
sum over substrate fields (BPS sector plus the
$\so(10)/\mathrm{SM}$ coset, all bosonic).  The matching
$\Lambda_{\mathrm{UV}} \approx \mPl$ is consistent with one
substrate mode per coset generator
($N_{\mathrm{DOF}}^{\mathrm{sub}} = 39$, cutoff at $0.983\,\mPl$);
exact equality ($N_{\mathrm{DOF}} = 12\pi$ exactly) is an
unattainable transcendental coincidence.  The Sakharov fermion-sign
issue of Table~\ref{tab:N-IR-counterexample} is an artifact of
treating IR-projected fermions as fundamental fields and does
\emph{not} apply at the substrate level.

This framing reinterprets the ``$N_{\mathrm{DOF}} = 12\pi$ matches
the GUT count of $39$ to $3.4\%$'' coincidence reported in earlier
drafts of this paper.  Earlier drafts asserted the matching
without distinguishing substrate from IR;\ the correct reading is
that the $39$ is a substrate-level generator count whose
identification as the right input to \eqref{eq:cutoff-shift}
awaits the substrate $c_J^{\mathrm{sub}}$ derivation.

\subsection{The structural identity $\mPl^2/\me^2 =
            |A_{K_7}|^{-2}$}
\label{sec:G3-K7-identity}

\begin{figure}[t]
\centering
\includegraphics[width=\textwidth]{paper18_fig_unified_torus.pdf}
\caption{The two substrate objects of the gravitational program
on the same Heegaard torus $T^2 \subset S^3 / 2I$.  (a) The $(2,3)$
torus knot (trefoil) is the BPS background of the matter Lagrangian
introduced in Paper~16:\ a $(p,q)$ winding on $T^2$ with mass
contribution $\propto p^2 + q^2$.  (b) The complete graph $K_7$ on
its Heffter triangular embedding (vertices labeled $0$--$6$,
$E = 21$ edges drawn as shortest-geodesic segments on the torus
surface) is the Wilson-graph object of Paper~17, evaluating to the
amplitude $A_{K_7} = \mathrm{bracket}(\alpha)\,\alpha^{21/2}$.  The
structural identity \eqref{eq:K7-identity},
$\mPl^2/\me^2 = |A_{K_7}|^{-2}$, says these two objects are not two
different substrate constructions but two views of one:\ the same
toroidal background that supports the BPS trefoil also supports
the $K_7$ graph state, and the relation between the matter-loop
cutoff (a) and the IR mass hierarchy (b) is fixed by their joint
embedding on $T^2$.}
\label{fig:K7-trefoil-on-torus}
\end{figure}

Independent of the cutoff question, the Phase~5 mismatch reported
in Paper~16 admits a clean structural reading.  Writing
Paper~17's closed form as a ratio:
\begin{equation}
  \frac{\mPl^2}{\me^2}
   \;=\;
   \frac{1}{\mathrm{bracket}(\alpha)^2\,\alpha^{21}}
   \;=\;
   \big|A_{K_7}\big|^{-2},
  \label{eq:K7-identity}
\end{equation}
where $A_{K_7} \equiv \mathrm{bracket}(\alpha)\,\alpha^{21/2}$ is
the Wilson amplitude evaluated by Paper~17 on the $K_7$ graph state
on the Heegaard torus of $S^3/2I$ (see
Figure~\ref{fig:K7-trefoil-on-torus} for the toroidal embedding of
both substrate objects).  Numerically,
\begin{equation}
  \big|A_{K_7}\big|^{-2}
   \;=\;
   137^{21}\,(8/7)^{-2}\,(1+\alpha/7+3\alpha^2)^{-2}
   \;\approx\;
   5.71 \times 10^{44},
  \label{eq:hierarchy-numeric}
\end{equation}
which agrees with the experimental Planck-to-electron mass-squared
hierarchy to $-29$\,ppm of CODATA.

\paragraph{What this is, and is not.}
Equation~\eqref{eq:K7-identity} is the Paper~17 result expressed in
a form that makes its identification with the Phase~5 hierarchy
manifest:\ the dimensionless number $5.71 \times 10^{44}$ \emph{is}
$\mPl^2/\me^2$ by definition, and Paper~17 derives it from the
$K_7$ graph state.  It is \emph{not} an independent derivation that the
Phase~5 gap \emph{must} equal $\big|A_{K_7}\big|^{-2}$;\ Paper~16's
Phase~5 reported a $\sim 10^{45}$ gap whose structural origin was
unidentified, and Paper~17 supplied that origin.  This paper takes
the structural identification as given and explores its dynamical
consequences (Sections~\ref{sec:G4-linearized}--\ref{sec:G6-einstein}).

The relevant question for the Sakharov program is whether the
\emph{matter loop} of the present paper is consistent with that
identification, i.e.\ whether the cutoff specified by
\eqref{eq:cutoff-shift} reproduces $\mPl$ to leading order.  As
established in Section~\ref{sec:G3-cutoff-shift}, this consistency
holds at the level of the $\sqrt{12\pi/N_{\mathrm{DOF}}}$ structural
relation, with a percent-level cutoff offset whose precise value
awaits the $c_J$-coefficient derivation noted there.

\subsection{Sketch of the corrected Phase~5 computation}
\label{sec:G3-phase5-sketch}

A satisfying confirmation that the $K_7$ amplitude is what
Paper~16's Phase~5 was missing is provided by re-doing the Phase~5
Casimir computation with the $K_7$ Lie-theoretic structure
inserted, and showing that the result lands on the experimentally
required $1/G$ rather than $10^{45}$ off.  We sketch the structural
form of that computation here; the full evaluation will appear
separately and the dimensional details below are schematic.

\paragraph{Phase~5 framework recap.}
Paper~16 \S Phase~5 evaluated the matter-loop induced $1/G$ on the
Brieskorn--Poincar\'e sphere $S^3/2I$ via the Casimir energy of
the BPS three-field theory, identifying the relevant background
as the trefoil $K_3$ at coupling $\lambda = e^2/2$.  Its scaling
form is
\begin{equation}
  \left[\frac{1}{16\pi G}\right]_{\mathrm{Phase\,5}}
   \;\propto\;
   \frac{C_2^{\mathrm{trefoil}}}{V_{S^3/2I}}\,\Lambda_{\mathrm{UV}}^2,
  \label{eq:phase5-form}
\end{equation}
where $C_2^{\mathrm{trefoil}}$ is the quadratic Casimir of the
$\SU(2)$-irrep carried by the trefoil-knot Wilson loop and
$V_{S^3/2I} = 2\pi^2/120$ is the volume of the Brieskorn sphere in
units where $R = 1$.  Numerically, with $\Lambda_{\mathrm{UV}} = \mPl$
identified naively, this evaluates to a number $\sim 10^{45}$ times
larger than the experimentally required $1/G$, the reported Phase~5
gap.

\paragraph{What Phase~5 was missing.}
The Casimir framework of Phase~5 captures the
\emph{quadratic-Casimir}, $\SU(2)$-trefoil contribution to the
matter loop on $S^3/2I$.  It does \emph{not} include:
\begin{itemize}
\item The Lie-theoretic structure from the full $\so(10)$
      embedding into $\mathrm{Spin}(7)$ that emerges in
      Paper~17's Wilson-amplitude computation.
\item The $K_7$ \emph{graph}-state structure linking the seven
      $\so(7)$ Cartan generators via cyclic exchange, which the
      Heegaard-torus computation of Paper~17 evaluates as
      $\mathrm{bracket}(\alpha)\,\alpha^{21/2}$.
\item The $\sin^{21}\!\bigl(\pi/(k+\hv)\bigr)$ Chern--Simons
      modular factor, which carries the cyclic-exchange
      information that the bare Casimir does not.
\end{itemize}
The Casimir computation lives at the trefoil-knot level; the
$K_7$ amplitude lives at the heptafoil-graph level.  The two
levels are separated by the $\mathrm{Spin}(7)$ representation theory
that Paper~17 unpacks, and we expect a missing dimensionless ratio
\begin{equation}
  \mathcal{R}_{K_7/\mathrm{trefoil}}
   \;\equiv\;
   \frac{|A_{K_7}|^{-2}}{C_2^{\mathrm{trefoil}}/V_{S^3/2I}}
  \label{eq:K7-trefoil-ratio}
\end{equation}
to enter the Phase~5 form when the $K_7$ Lie-theoretic structure is
restored.  Inserting it converts \eqref{eq:phase5-form} into the
IR-measurable form of Paper~17, schematically
\begin{equation}
  \left[\frac{1}{16\pi G}\right]_{\mathrm{Phase\,5}}^{\!\!\!\mathrm{corrected}}
   \;\sim\;
   \mathcal{R}_{K_7/\mathrm{trefoil}}\,\cdot\,
   \frac{C_2^{\mathrm{trefoil}}}{V_{S^3/2I}}\,\mPl^2
   \;=\;
   |A_{K_7}|^{-2}\,\frac{\mPl^2}{V_{S^3/2I}},
  \label{eq:phase5-corrected}
\end{equation}
with $|A_{K_7}|^{-2} = \mPl^2/\me^2$ supplying the $10^{45}$ factor
that the bare Casimir framework was missing.

\paragraph{Status.}
The full Phase~5 corrected computation requires evaluating
$\mathcal{R}_{K_7/\mathrm{trefoil}}$ from first principles in the
Chern--Simons functional-integral formulation on $S^3/2I$,
following the construction outlined in Paper~17 \S6 but applied
to the matter-loop coefficient rather than the
$\me/\mPl$-fixing observable.  We do not perform that computation
in this paper; it is a separate technical exercise that would
strengthen the structural picture established here but is not
required for the dynamical content of
Sections~\ref{sec:G4-linearized}--\ref{sec:G6-einstein}.  The
corrected Phase~5 computation is logged as a follow-up.

\subsection{Numerical comparison: $G$ to $+29$\,ppm of CODATA}

We can now compute Newton's constant numerically from the unified
expression.  Writing
\begin{equation}
  G
   \;=\;
   \frac{\mathrm{bracket}(\alpha)^2\,\alpha^{21}}{\me^2}\,(\hbar c),
  \label{eq:G-final}
\end{equation}
and using CODATA $\alpha = 7.2973525693\times10^{-3}$,
$\me = 0.51099895\,\mathrm{MeV}/c^2$, the
\texttt{nwt\_substrate.gravity.}\allowbreak\texttt{G\_\allowbreak substrate\_\allowbreak SI()}
routine evaluates
\begin{equation}
  G_{\mathrm{NWT,\,NNLO}}
   \;=\;
   6.674494 \times 10^{-11}\;\mathrm{m^3\,kg^{-1}\,s^{-2}}.
\end{equation}
Compared with the CODATA experimental value
$G_{\mathrm{CODATA}} = 6.67430(15) \times 10^{-11}$ m$^3$/kg/s$^2$:
\begin{equation}
  \frac{G_{\mathrm{NWT,\,NNLO}} - G_{\mathrm{CODATA}}}{G_{\mathrm{CODATA}}}
   \;=\;
   +29\,\mathrm{ppm}.
\end{equation}
The $\pm 22$\,ppm experimental uncertainty on $G$ (the largest
relative uncertainty among the fundamental constants) places the
NWT prediction \emph{inside} the experimental error bar.  We
emphasize that this numerical agreement is a re-expression of
Paper~17's $\me/\mPl$ closed form via $G\,\me^2 =
\mathrm{bracket}^2\,\alpha^{21}$;\ what the Sakharov derivation of
this paper adds is a \emph{dynamical} effective action consistent
with the Paper~17 number at leading order, not an independent
numerical re-derivation of $G$.

\subsection{Summary}

The structural picture established in this section is:
\begin{itemize}
\item The Sakharov matter loop produces the long-wavelength
      Einstein--Hilbert action with coefficient
      $1/(16\pi G) = N_{\mathrm{DOF}}\Lambda_{\mathrm{UV}}^2 / (192\pi^2)$,
      with $N_{\mathrm{DOF}}$ the $c_J$-weighted matter content.
\item Paper~17's $K_7$ Wilson amplitude
      $\mathrm{bracket}(\alpha)\alpha^{21/2}$ specifies $\me/\mPl$ in
      IR-measurable terms; equivalently
      $\mPl^2/\me^2 = |A_{K_7}|^{-2}$.
\item Equating the two routes gives the structural relation
      $\Lambda_{\mathrm{UV}}/\mPl =
      \sqrt{12\pi/N_{\mathrm{DOF}}^{\mathrm{sub}}}$, with
      $N_{\mathrm{DOF}}^{\mathrm{sub}}$ a $c_J$-weighted sum over
      \emph{substrate} matter fields (not IR projections).
      One-substrate-mode-per-coset-generator counting gives
      $N_{\mathrm{DOF}}^{\mathrm{sub}} = 39$ and
      $\Lambda_{\mathrm{UV}} = 0.983\,\mPl$;\ a substrate-level
      heat-kernel derivation of $c_J^{\mathrm{sub}}$ is logged as a
      follow-up (Section~\ref{sec:G3-cutoff-shift}).
\item The $10^{45}$ Phase~5 gap of Paper~16 is structurally
      identified with $|A_{K_7}|^{-2}$, the missing IR
      Lie-theoretic factor that the pure $S^3/2I$ Casimir framework
      did not include.  A corrected Phase~5 computation
      incorporating this factor is sketched in
      Section~\ref{sec:G3-phase5-sketch}; the full evaluation is a
      logged follow-up.
\item Numerically, $G_{\mathrm{NWT,\,NNLO}} = +29$\,ppm of CODATA
      (inside the $\pm 22$\,ppm experimental error), consistent
      with Paper~17's $\me/\mPl$ closed form.
\end{itemize}

We have thus produced the gravitational coupling at the
\emph{coefficient} level.  The remaining Sakharov program---the
linearized Einstein equation, the gravitational-wave dispersion,
the curved-background extension, and the full nonlinear Einstein
equations as a variational consequence of the induced
$\Gamma[g]$---is carried out in the following sections.

% ==========================================================
\section{Linearised Einstein equations and gravitational waves (G4)}
\label{sec:G4-linearized}
% ==========================================================

With the EH coefficient $1/(16\pi G)$ fixed by the
matter-loop / $K_7$-amplitude bridge of
Section~\ref{sec:G3-K7-bridge}, the Sakharov-induced effective
action contains, at single-curvature order,
\begin{equation}
  \Gamma_{\mathrm{EH}}[g]
   \;=\;
   -\,\frac{1}{16\pi G}\int d^4x\,\sqrt{-g}\,R(g).
  \label{eq:gamma-EH-recap}
\end{equation}
Linearising around flat space, $g_{\mu\nu} = \eta_{\mu\nu} +
h_{\mu\nu}$, yields the standard Fierz--Pauli kinetic action for
the metric perturbation, the linearized Einstein equation, the
gravitational-wave dispersion, and Newton's law in the
non-relativistic static limit.  The detailed checks are in
\texttt{analysis/}\allowbreak\texttt{nwt\_\allowbreak paper18\_\allowbreak G4\_\allowbreak linearized\_\allowbreak einstein.py};
we give the analytic content here.

\subsection{Linearised curvature and the Fierz--Pauli action}

To first order in $h_{\mu\nu}$, the Ricci scalar reads
\begin{equation}
  R^{(1)}
   \;=\;
   \partial_\mu\partial_\nu\,h^{\mu\nu}
   \;-\;
   \Box\,h,
   \qquad
   h \equiv \eta^{\mu\nu}\,h_{\mu\nu}.
  \label{eq:R-linear}
\end{equation}
Inserting \eqref{eq:R-linear} into \eqref{eq:gamma-EH-recap} and
integrating by parts, the linearized Einstein--Hilbert action
becomes the Fierz--Pauli kinetic action,
\begin{equation}
  S_{\mathrm{FP}}
   \;=\;
   \frac{1}{64\pi G}
   \int d^4 x\,\Bigl[\,
     -\tfrac{1}{2}\,\partial_\alpha h_{\mu\nu}\,\partial^\alpha h^{\mu\nu}
     + \partial^\mu h_{\mu\nu}\,\partial^\alpha {h^\nu}{}_\alpha
     - \partial^\mu h_{\mu\nu}\,\partial^\nu h
     + \tfrac{1}{2}\,\partial^\alpha h\,\partial_\alpha h
  \,\Bigr].
  \label{eq:Fierz-Pauli}
\end{equation}
This is the unique Lorentz-invariant 2-tensor action that
propagates exactly the two transverse-traceless polarizations of a
massless graviton; any other symmetric-2-tensor kinetic structure
would propagate ghost or tachyon modes.  The Sakharov mechanism
generates the Fierz--Pauli structure automatically because the
matter sector is unitary (Section~\ref{sec:G2-flat-sakharov}); we
therefore inherit ghost-freedom from the matter sector itself.

\subsection{The linearized Einstein equation}

Adding the matter--graviton coupling \eqref{eq:Smatter-linear} and
varying the total action $S_{\mathrm{FP}} +
S_{\mathrm{matter}}^{(1)}$ with respect to $h^{\mu\nu}$:
\begin{equation}
  \frac{1}{16\pi G}\,G^{(1)}_{\mu\nu}[h]
   \;=\;
   \tfrac{1}{2}\,T^{(M)}_{\mu\nu},
  \label{eq:linearized-EE}
\end{equation}
where the linearized Einstein tensor is
\begin{align}
  G^{(1)}_{\mu\nu}[h]
   \;=\; -\tfrac{1}{2}\,\Big[\,
     &\Box h_{\mu\nu}
     - \partial_\mu\partial^\alpha h_{\alpha\nu}
     - \partial_\nu\partial^\alpha h_{\alpha\mu} \nonumber\\
     & + \partial_\mu\partial_\nu h
     - \eta_{\mu\nu}\,\bigl(\Box h - \partial^\alpha\partial^\beta h_{\alpha\beta}\bigr)
   \Big].
   \label{eq:G-linear}
\end{align}
Equation~\eqref{eq:linearized-EE} is the linearized Einstein
equation with the $G$ of \eqref{eq:G-final}.  The factor $1/2$
on the matter side is the standard Sakharov normalization
matching the $-\tfrac{1}{2}\,h^{\mu\nu} T^{(M)}_{\mu\nu}$ vertex
from G1.

\subsection{Transverse-traceless gauge and gravitational waves}

Diffeomorphism invariance of \eqref{eq:linearized-EE} permits the
gauge transformation $h_{\mu\nu} \to h_{\mu\nu}
+ \partial_\mu \xi_\nu + \partial_\nu \xi_\mu$.  Choosing $\xi^\mu$
to enforce de~Donder gauge $\partial^\mu h_{\mu\nu} = \tfrac{1}{2}
\partial_\nu h$, plus the residual freedom to set
$\eta^{\mu\nu} h_{\mu\nu} = 0$, leaves only the
transverse-traceless components,
\begin{equation}
  \partial^\mu h^{\mathrm{TT}}_{\mu\nu} = 0,
  \qquad
  \eta^{\mu\nu} h^{\mathrm{TT}}_{\mu\nu} = 0.
  \label{eq:TT-gauge}
\end{equation}
Counting: a symmetric $4\times 4$ tensor has $10$ components;
diffeomorphism invariance removes $4$, the residual harmonic gauge
removes another $4$ (one for each component of $\xi^\mu$ satisfying
$\Box \xi^\mu = 0$).  The result, $10 - 4 - 4 = 2$, is the
standard count of two physical TT polarizations of a massless
graviton---one each for $h_+$ and $h_\times$.

In TT gauge, equation~\eqref{eq:linearized-EE} collapses to the
gravitational-wave equation
\begin{equation}
  \boxed{\;
    \Box\,h^{\mathrm{TT}}_{\mu\nu}
     \;=\;
     -16\pi G\,T^{\mathrm{TT}}_{\mu\nu}.
  \;}
  \label{eq:GW-equation}
\end{equation}
On-shell graviton modes ($T^{\mathrm{TT}} = 0$) satisfy
$p_\mu p^\mu = 0$, the masslessness condition.  The two TT
polarizations propagate at the speed of light.  This is the
standard linearized Einstein equation, recovered with the
structurally derived $G$ of \eqref{eq:G-final}.

\subsection{Newton's law in the non-relativistic limit}

For a static, weak-field source with stress-energy dominated by a
mass density $\rho$, the linearized Einstein equation reduces in
the Newtonian limit to
\begin{equation}
  \nabla^2\,\Phi(\vec x)
   \;=\;
   4\pi G\,\rho(\vec x),
  \label{eq:Newton}
\end{equation}
where $\Phi$ is the Newtonian potential identified by the
$00$-component of the metric perturbation,
$h_{00} = -2\Phi$.  For two point masses $m_1, m_2$ at separation
$r$, the resulting force is
\begin{equation}
  F(r)
   \;=\;
   -\,\frac{G\,m_1 m_2}{r^2},
\end{equation}
with the structurally derived $G$ of \eqref{eq:G-final}.  The
classical $1/r^2$ universal-attraction law of Newton, the
gravitational red-shift of static observers, and all standard
post-Newtonian phenomenology follow from \eqref{eq:Newton} as in
the textbook treatment.

\subsection{Summary}

The linearized dynamics established in this section---Fierz--Pauli
kinetic action, linearized Einstein equation, TT gravitational
waves with two physical polarizations, Newtonian limit---are the
standard results of weak-field general relativity, recovered here
with the $G$ derived structurally from the $K_7$ Wilson amplitude.
What this section adds beyond textbook GR is the chain of
connections from the $K_7$ amplitude to gravitational radiation:
Paper~17's $K_7$ Wilson product on the Heegaard torus of $S^3/2I$
specifies $\me/\mPl$, the matter loop produces the Fierz--Pauli
kinetic term with the $1/G$ coefficient implied by the consistency
condition of Section~\ref{sec:G3-K7-bridge}, and the linearized
Einstein equation \eqref{eq:linearized-EE} then governs the
propagation of small $h_{\mu\nu}$ perturbations on $\eta_{\mu\nu}$.
The next step is to extend this to a general curved background and
pick up the higher-curvature corrections.

% ==========================================================
\section{Sakharov on a curved background:
higher-curvature and the cosmological constant (G5)}
\label{sec:G5-curved}
% ==========================================================

The flat-space derivation of Sections~\ref{sec:G2-flat-sakharov}
and~\ref{sec:G3-K7-bridge} fixed the coefficient of $\sqrt{-g}\,R$
in the Sakharov-induced effective action.  Extending the matter
loop to an arbitrary curved background $g_{\mu\nu}$ produces, in
addition to the linear-curvature term, a cosmological constant
$\Lambda_{\mathrm{cc}}$ and Stelle-type higher-curvature
corrections~\cite{stelle1977}.  These additional terms arise
naturally from the heat-kernel expansion on curved space; we
catalogue them here following the Vassilevich
manual~\cite{vassilevich2003}.  The full computation lives in
\texttt{analysis/nwt\_paper18\_G5\_curved\_sakharov.py}.

\subsection{The heat-kernel expansion in curved space}

For a Laplace-type operator $K = -\Box + E(x)$ on a 4-manifold $M$,
where $E$ is the bundle endomorphism that collects mass terms and
curvature couplings, Vassilevich's general result
\cite{vassilevich2003} gives the Seeley--DeWitt coefficients
\begin{align}
  a_0(K) &= \mathrm{tr}\,\mathbb{1}, \\
  a_1(K) &= \tfrac{R}{6}\,\mathrm{tr}\,\mathbb{1} - \mathrm{tr}\,E, \\
  a_2(K) &= \tfrac{1}{180}\,\mathrm{tr}\,\mathbb{1}\,
            \bigl(R^2 + R_{\mu\nu\rho\sigma}^2 - R_{\mu\nu}^2\bigr)
          + \tfrac{1}{2}\,\mathrm{tr}\bigl(E - \tfrac{R}{6}\mathbb{1}\bigr)^2
          + \text{boundary terms}.
\end{align}
The flat-space limit recovers the Section~\ref{sec:G2-flat-sakharov}
expressions (the $a_1$ term reduces to $-m^2$ on flat space, and
$a_2$ collapses to $\tfrac{1}{2}m^4$ plus boundary terms).  In
curved space, the additional curvature-squared structures in $a_2$
generate Stelle higher-curvature corrections, and the
$\mathrm{tr}\,\mathbb{1}$ piece of $a_0$ generates a cosmological
constant.

Inserting these coefficients into the cutoff-regulated effective
action \eqref{eq:gamma-cutoff}, the resulting curvature-dependent
pieces of $\Gamma[g]$ are
\begin{equation}
\Gamma_{\mathrm{grav}}[g]
   \;=\;
   -\int d^4 x\,\sqrt{-g}\,
   \Bigl[\,
       \Lambda_{\mathrm{cc}}
     + \tfrac{1}{16\pi G}\,R(g)
     + \alpha_S\,R^2
     + \beta_S\,R_{\mu\nu} R^{\mu\nu}
     + \gamma_S\,R_{\mu\nu\rho\sigma} R^{\mu\nu\rho\sigma}
     + \cdots
   \Bigr],
  \label{eq:gamma-curved}
\end{equation}
where the coefficients $\Lambda_{\mathrm{cc}}$, $1/G$,
$\alpha_S$, $\beta_S$, $\gamma_S$ depend on the matter content and
the cutoff.

\subsection{The Einstein--Hilbert coefficient is matter-independent
in the curved extension}

By comparison of \eqref{eq:gamma-curved} with the flat-space
$a_1$-coefficient extraction of \eqref{eq:per-scalar}, the linear
$R$-coefficient is exactly the flat-space result of
Section~\ref{sec:G2-flat-sakharov}:
\begin{equation}
  \frac{1}{16\pi G}
   \;=\;
   \frac{N_{\mathrm{DOF}}\,\Lambda_{\mathrm{UV}}^2}{192\pi^2}.
  \label{eq:G-curved}
\end{equation}
The flat-space and curved-space Sakharov calculations are
self-consistent at the EH level, with the same $G$ obtained in
either case.  The cutoff identified through Paper~17's $K_7$
amplitude (Section~\ref{sec:G3-K7-bridge}) carries through
unchanged.

\subsection{Stelle higher-curvature corrections}

The $a_2$ coefficient generates the higher-curvature pieces.
Their precise coefficients depend on the spin content; for the
NWT matter sector,
\begin{align}
  \alpha_S &= -\frac{N_{\mathrm{DOF}}}{180 \cdot 32\pi^2}\,
              \ln\!\bigl(\Lambda_{\mathrm{UV}}^2/m^2\bigr) + \cdots, \\
  \beta_S, \gamma_S &= \text{analogous log-divergent coefficients
                                 from $R_{\mu\nu}^2$, $R_{\mu\nu\rho\sigma}^2$ pieces}.
\end{align}
At low energies $R \ll \Lambda_{\mathrm{UV}}^2 = \mPl^2$ the
higher-curvature contributions are subdominant to the
Einstein--Hilbert term by powers of $R/\mPl^2 \sim 10^{-44}$ at
laboratory scales.  Stelle showed in 1977~\cite{stelle1977} that
the $R^2 + R_{\mu\nu}^2$ structure is renormalisable but introduces
ghost modes near the Planck scale.  The ghost-mode issue is
generic to higher-curvature gravity and is not specific to NWT;
we flag it as an open problem in Section~\ref{sec:open} and
proceed with the Einstein--Hilbert truncation as physically
relevant for all $R \ll \mPl^2$.

\subsection{The cosmological-constant coefficient}

The $a_0$ coefficient generates the matter-loop vacuum-energy
contribution to the cosmological constant,
\begin{equation}
  \Lambda_{\mathrm{cc}}^{\mathrm{Sakharov}}
   \;=\;
   \frac{N_{\mathrm{DOF}}\,\Lambda_{\mathrm{UV}}^4}{64\pi^2}.
  \label{eq:Lambda-cc-Sakharov}
\end{equation}
With $\Lambda_{\mathrm{UV}} = \mPl$ as fixed by
Section~\ref{sec:G3-K7-bridge} and $N_{\mathrm{DOF}} = 12\pi$,
this gives $\Lambda_{\mathrm{cc}}^{\mathrm{Sakharov}}
\sim \mPl^4 \approx 2 \times 10^{72}\,\mathrm{GeV}^4$,
which is approximately $10^{120}$ times larger than the observed
cosmological constant
$\Lambda_{\mathrm{cc}}^{\mathrm{obs}} \approx
3 \times 10^{-47}\,\mathrm{GeV}^4$.

This is the standard cosmological-constant problem (the $\Lambda
\sim \mPl^4$ vacuum-energy mismatch).  It is \emph{not} specific
to NWT: any Sakharov-induced gravity construction with
$\Lambda_{\mathrm{UV}} = \mPl$ produces $\Lambda_{\mathrm{cc}}
\sim \mPl^4$.  Resolving it requires substrate-level cancellations
(Coleman--Weinberg-type partial closures of the vacuum-energy
hierarchy on the $K_7$ structure are explored in Paper~21,
deferred from this paper).  We acknowledge the gap, log it as
unresolved (Section~\ref{sec:open}), and note that the
cosmological-constant problem is a problem of Sakharov-induced
gravity in general, not of the NWT specification of it.

The cosmological constant is also \emph{independent} of the
Einstein--Hilbert coefficient: the $a_0$ coefficient
\eqref{eq:Lambda-cc-Sakharov} and the $a_1$ coefficient
\eqref{eq:G-curved} are decoupled in the Vassilevich formalism, so
the structural derivation of $G$ in
Sections~\ref{sec:G2-flat-sakharov}--\ref{sec:G3-K7-bridge} is not
contaminated by the unresolved $\Lambda_{\mathrm{cc}}$ problem.

\subsection{Summary}

Extending the Sakharov calculation to a generic curved background
produces an effective action containing the Einstein--Hilbert
term, Stelle higher-curvature corrections, and a cosmological
constant.  The Einstein--Hilbert coefficient is the same as on
flat space: \eqref{eq:G-curved} agrees with
Section~\ref{sec:G2-flat-sakharov}, and the $K_7$ amplitude bridge
of Section~\ref{sec:G3-K7-bridge} carries through unchanged.  The
higher-curvature corrections and cosmological constant are
generated in addition; the former are subdominant for
$R \ll \mPl^2$, the latter is the standard
$\Lambda_{\mathrm{cc}} \sim \mPl^4$ problem deferred to Paper~21.
With the curved-space effective action
\eqref{eq:gamma-curved} now in hand, we proceed to derive the
nonlinear Einstein equations as the variational consequence.

% ==========================================================
\section{Variation of the effective action: full nonlinear
Einstein equations (G6)}
\label{sec:G6-einstein}
% ==========================================================

The Sakharov-induced gravitational action established in
Sections~\ref{sec:G2-flat-sakharov}--\ref{sec:G5-curved} is, in
its truncation to single-curvature order,
\begin{equation}
  \Gamma_{\mathrm{grav}}[g]
   \;=\;
   -\int d^4 x\,\sqrt{-g}\,
   \Bigl[\,
       \Lambda_{\mathrm{cc}}
     + \frac{1}{16\pi G}\,R(g)
   \Bigr]
   \;+\;\Gamma_{\mathrm{HC}}[g],
  \label{eq:gamma-truncated}
\end{equation}
where $\Gamma_{\mathrm{HC}}$ collects the Stelle higher-curvature
contributions, subdominant for $R \ll \mPl^2$.  The full nonlinear
Einstein equations follow from the variational principle
$\delta(\Gamma_{\mathrm{grav}} + S_{\mathrm{matter}})/\delta g^{\mu\nu}
= 0$ applied to \eqref{eq:gamma-truncated}.  This is textbook
general relativity \emph{except} that the gravitational coupling
$1/(16\pi G)$ is now structurally derived rather than postulated.
The full computation lives in
\texttt{analysis/}\allowbreak\texttt{nwt\_\allowbreak paper18\_\allowbreak G6\_\allowbreak einstein\_\allowbreak variation.py}.

\subsection{Variation of the effective action}

The Einstein--Hilbert variation proceeds standardly:
\begin{equation}
  \delta\bigl(\sqrt{-g}\,R\bigr)
   \;=\;
   \sqrt{-g}\,\bigl(R_{\mu\nu} - \tfrac{1}{2}\,g_{\mu\nu} R\bigr)\,\delta g^{\mu\nu}
   \;+\;
   \sqrt{-g}\,\bigl(g^{\mu\nu}\Box\,\delta g_{\mu\nu}
                     - \nabla^\mu\nabla^\nu\,\delta g_{\mu\nu}\bigr).
\end{equation}
The second term is a total derivative and contributes only the
boundary Gibbons--Hawking--York term, which we drop for closed
or asymptotically-flat manifolds.  The cosmological-constant
piece varies as
$\delta(\sqrt{-g}\,\Lambda_{\mathrm{cc}}) =
-\tfrac{1}{2}\,\sqrt{-g}\,\Lambda_{\mathrm{cc}}\,g_{\mu\nu}\,
\delta g^{\mu\nu}$, and the matter piece by definition contributes
$\delta S_{\mathrm{matter}}/\delta g^{\mu\nu}
= -\tfrac{1}{2}\,\sqrt{-g}\,T_{\mu\nu}$.

Setting $\delta(\Gamma_{\mathrm{grav}} + S_{\mathrm{matter}})/\delta g^{\mu\nu}
= 0$ and solving:
\begin{equation}
  \boxed{\;
    G_{\mu\nu}
     \;+\;
     \Lambda_{\mathrm{cc}}\,g_{\mu\nu}
     \;=\;
     8\pi G\,T_{\mu\nu}
     \;+\;
     \alpha_S\,\partial_\mu\partial_\nu R + \cdots,
  \;}
  \label{eq:Einstein-full}
\end{equation}
the full nonlinear Einstein equations with the
structurally-derived $G$ of \eqref{eq:G-final}, the cosmological
constant of \eqref{eq:Lambda-cc-Sakharov}, and the Stelle
higher-curvature corrections from $\Gamma_{\mathrm{HC}}$.  This is
the central dynamical statement of the paper.

The non-linearity is purely geometric: $G_{\mu\nu} \equiv R_{\mu\nu}
- \tfrac{1}{2}\,g_{\mu\nu}\,R$ depends nonlinearly on $g_{\mu\nu}$
through the Riemann tensor.  No additional matter physics is
required.  Linearisation $g = \eta + h$ in
\eqref{eq:Einstein-full} reproduces the linearized result
\eqref{eq:linearized-EE} of Section~\ref{sec:G4-linearized}, with
$\Lambda_{\mathrm{cc}}$ contributing a constant term that we set
to zero in the laboratory limit.

\subsection{Bianchi identity and conservation of $T_{\mu\nu}$}

The contracted Bianchi identity
\begin{equation}
  \nabla^\mu G_{\mu\nu} \;=\; 0
\end{equation}
is a geometric statement valid on any metric; combined with
$\nabla^\mu(\Lambda_{\mathrm{cc}}\,g_{\mu\nu}) = 0$ (since
$\Lambda_{\mathrm{cc}}$ is a constant and $\nabla g = 0$), it
implies
\begin{equation}
  \nabla^\mu T_{\mu\nu} \;=\; 0,
\end{equation}
the local conservation of matter energy-momentum.  This holds
nonlinearly, automatically, as a consistency consequence of the
Einstein equations rather than an additional postulate.  It also
provides a non-trivial check on the matter sector
(Section~\ref{sec:matter-lagrangian}): the $T^{(M)}_{\mu\nu}$ of
the BPS L1+L2 fluctuations must satisfy $\nabla^\mu
T^{(M)}_{\mu\nu} = 0$ on its equations of motion, which it does
identically by virtue of the gauge invariance of the matter
Lagrangian.

\subsection{Schwarzschild as a vacuum solution: symbolic
verification}

The Schwarzschild metric
\begin{equation}
  ds^2
   \;=\;
   -\bigl(1 - \tfrac{2GM}{r}\bigr)\,dt^2
   + \bigl(1 - \tfrac{2GM}{r}\bigr)^{-1}\,dr^2
   + r^2\,d\Omega^2,
  \label{eq:Schwarzschild}
\end{equation}
is the textbook static, spherically symmetric vacuum solution.  We
verify it as a vacuum solution of \eqref{eq:Einstein-full} (with
$\Lambda_{\mathrm{cc}} = 0$ and $\Gamma_{\mathrm{HC}} = 0$) by
direct symbolic computation.  Calling
\texttt{nwt\_substrate.gravity.}\allowbreak\texttt{verify\_\allowbreak schwarzschild\_\allowbreak vacuum\_\allowbreak symbolic()}
performs the full sympy derivation:
\begin{itemize}
\item The metric tensor and its inverse are constructed in
      $(t, r, \theta, \phi)$ coordinates.
\item All $40$ Christoffel symbols
      $\Gamma^{\rho}{}_{\mu\nu} = \tfrac{1}{2}\,g^{\rho\sigma}
      (\partial_\mu g_{\nu\sigma} + \partial_\nu g_{\mu\sigma}
      - \partial_\sigma g_{\mu\nu})$ are computed; their
      symmetric and metric-compatible properties are verified
      symbolically.
\item The $20$ independent Riemann tensor components
      $R^{\rho}{}_{\sigma\mu\nu}$ are evaluated.
\item The Ricci tensor $R_{\mu\nu} = R^\rho{}_{\mu\rho\nu}$ is
      contracted; all $10$ independent components of the symmetric
      tensor vanish identically:
      $R_{\mu\nu}(g_{\mathrm{Sch}}) = 0$ for all
      $\mu, \nu \in \{0, 1, 2, 3\}$.
\end{itemize}
The verification confirms that NWT's Sakharov-induced gravitational
sector admits Schwarzschild as a vacuum solution.  Under the
identification $r_s = 2 G M / c^2$, the Schwarzschild radius for
any mass $M$ is governed by the structurally-derived $G$ of
\eqref{eq:G-final}.

\paragraph{Sanity check: NWT particles are not black holes.}  In
Planck units, $r_s = 2 M$ for any mass $M$.  The Schwarzschild
radius of an electron is
$r_s^{(e)} = 2\me / \mPl^2 \approx 1.4 \times 10^{-57}\,\mathrm{m}$,
many orders of magnitude smaller than the electron Compton
wavelength $\xi_e = \hbar/(\me c) \approx 3.9 \times 10^{-13}\,\mathrm{m}$
and the corresponding torus-knot vortex correlation length
(Paper~6).  Gravitational backreaction is negligible at
fundamental-particle scales.  Astrophysical-mass black holes
($M \gtrsim M_\odot$) live in the regime $r_s \sim$ km where
\eqref{eq:Einstein-full} reduces to standard general relativity
with the structurally derived $G$.  No fundamental NWT particle
is a black hole.

\subsection{Summary}

Variation of the Sakharov-induced effective action with respect to
$g^{\mu\nu}$ yields the full nonlinear Einstein equations
\eqref{eq:Einstein-full}, with the structurally-derived $G$ of
\eqref{eq:G-final} and the cosmological constant of
\eqref{eq:Lambda-cc-Sakharov}; matter energy-momentum conservation
$\nabla^\mu T_{\mu\nu} = 0$ follows as a consistency consequence
of the contracted Bianchi identity rather than as an additional
postulate.  The Schwarzschild metric is verified as a vacuum
solution by symbolic computation, all $10$ independent components
of the Ricci tensor vanishing identically.  The four-paper NWT
gravitational program (Papers~14--18) extends the structural
Paper~14--17 picture to a dynamical theory: gravity is identified
with the long-wavelength variational dynamics of the
Sakharov-induced effective action, with the open problems of
Section~\ref{sec:open}---cosmological-constant fine-tuning,
ghost-freedom of the Stelle higher-curvature sector beyond linear
order, and a quantum-level treatment of diffeomorphism invariance
on curved backgrounds---left for follow-up work.

% ==========================================================
\section{Open problems}
\label{sec:open}
% ==========================================================

Four open problems flagged in the course of the derivation, none
of which are NWT-specific but all of which interact with the
gravitational program assembled here.

\subsection{The cosmological-constant fine-tuning}

The Sakharov mechanism on a curved background generates the matter
loop's vacuum-energy contribution to the cosmological constant
\eqref{eq:Lambda-cc-Sakharov}.  With the cutoff
$\Lambda_{\mathrm{UV}} \approx \mPl$ and
$N_{\mathrm{DOF}}^{\mathrm{sub}} \sim \mathcal{O}(10)$, this gives
$\Lambda_{\mathrm{cc}} \sim \mPl^4$, approximately $10^{120}$ times
the observed value
$\Lambda_{\mathrm{cc}}^{\mathrm{obs}} \approx 3 \times 10^{-47}$
GeV$^4$.  This is the standard cosmological-constant problem of
Sakharov-induced gravity, not specific to NWT.

The $K_7$ amplitude that bridges UV and IR for $1/G$ supplies a
factor of $|A_{K_7}|^{-2} \approx 5.7 \times 10^{44}$, which would
nominally close $\sim 45$ orders of magnitude of the
$\sim 10^{120}$ gap if a Coleman--Weinberg-type cancellation in
the vacuum-energy sector is mediated by an analogous structural
factor (perhaps as $\mPl^{4-2k}\,\me^{2k}$ for some integer $k$).
The remaining $\sim 75$ orders are not addressed by any structure
established in the present paper.  A serious treatment of the
$\Lambda_{\mathrm{cc}}$ problem within the substrate-monism
framework is the subject of Paper~21~\cite{paper19} and we do not
attempt resolution here;\ the order-of-magnitude framing above is
suggestive of a structural mechanism rather than a partial
solution.

\subsection{Black hole formation from vortex collapse}

The Schwarzschild verification (Section~\ref{sec:G6-einstein})
established that NWT's gravitational sector admits black holes as
\emph{vacuum} solutions.  An open question is whether dynamical
collapse of an NWT vortex configuration to a black hole is
mechanistically describable within the Sakharov-induced framework.
Two qualitative observations bound the question:
no fundamental NWT particle is itself a black hole (the
Schwarzschild radius of an electron is $\sim 10^{-57}\,\mathrm{m}$,
far below the electron Compton scale); and stellar-mass
collapses, dominated at $r_s \sim$ km by the long-wavelength
limit of \eqref{eq:Einstein-full}, must reproduce standard general
relativity.  The intermediate regime---a substrate-resolved
description of the collapse process itself---is open and is the
motivation for the cosmogenesis program of Paper~22.

\subsection{Stelle higher-curvature ghost-freedom beyond linear order}

The Stelle higher-curvature corrections \eqref{eq:gamma-curved}
introduce $R^2$, $R_{\mu\nu}^2$, $R_{\mu\nu\rho\sigma}^2$ kinetic
structures that are renormalisable~\cite{stelle1977} but propagate
ghost modes near the Planck scale.  For all $R \ll \mPl^2$ the
higher-curvature sector is suppressed by powers of $R/\mPl^2$ and
the Einstein--Hilbert truncation is physically valid; the ghost
modes are an issue only at curvatures of Planck-scale order and
above.  This is a known issue in higher-derivative gravity in
general and is not specific to NWT.

\paragraph{NWT-specific stance.}
Within the framework of the present paper, three observations
constrain the status of the Stelle ghosts:
\begin{itemize}
\item \textbf{The Stelle Lagrangian is an EFT truncation of the
      Sakharov effective action, not a fundamental theory.}
      The $R^2$ and Weyl-squared terms are heat-kernel $a_2$
      coefficients with finite, calculable prefactors set by the
      matter spectrum (Sections~\ref{sec:G5-curved}); they are not
      independent gravitational degrees of freedom in the
      microscopic substrate Lagrangian of Paper~16.  The ghost modes
      are therefore artifacts of the long-wavelength expansion past
      its regime of validity, not physical degrees of freedom of
      the substrate.
\item \textbf{The substrate's matter spectrum is bosonic at BPS
      coupling.}  The Bogomolny condition $\lambda = e^2/2$ that
      defines the matter Lagrangian saturates a stability bound;\
      the BPS structure constrains the form of the
      higher-curvature coefficients (in particular, the ratio
      $\alpha_S/\alpha_W$ between the $R^2$ and Weyl-squared
      coefficients depends on the BPS spectrum and is not freely
      tunable).  Whether the BPS structure forces a specific sign or
      cancellation pattern that suppresses the ghosts at Planck-scale
      curvature is an open question;\ at present we have only the
      observation that the BPS coefficients are constrained, not the
      derivation that they are constrained in the right direction.
\item \textbf{The Sakharov derivation breaks down at curvatures
      $R \sim \mPl^2$ for an independent reason.}  At such curvatures
      the heat-kernel proper-time expansion no longer converges and
      the IR-projected matter spectrum is no longer the right input
      to the heat-kernel sum;\ the substrate-projection mechanism
      itself becomes the relevant physics.  In other words, the EFT
      regime ends before the Stelle ghosts become physically
      observable, and any consistent treatment of the substrate
      dynamics at $R \sim \mPl^2$ would replace the Stelle Lagrangian
      with a substrate-level computation that we have not performed.
\end{itemize}
The conservative reading is therefore that the Stelle ghosts are an
EFT artifact rather than a physical instability, and that resolution
will come from the substrate-level dynamics replacing the EFT at
high curvature rather than from a ghost-freedom theorem proved within
the EFT itself.  Possible literature resolutions (Stelle's asymptotically-free fixed
point, ghost-resummation in non-local-gravity proposals) apply
equally to the NWT-specific Sakharov sector and may be imported as
needed for specific phenomenological purposes;\ the present paper
does not commit to any one of them.

\subsection{Background independence and quantum diffeomorphism invariance}

The Sakharov derivation of Sections~\ref{sec:G2-flat-sakharov}
through~\ref{sec:G6-einstein} treats the metric $g_{\mu\nu}$ as a
classical background field on which the matter loop is evaluated;\
the metric perturbation $h_{\mu\nu}$ is then introduced as a
classical fluctuation around flat (Section~\ref{sec:G4-linearized})
or arbitrary curved (Section~\ref{sec:G5-curved}) background.  The
emergent Einstein--Hilbert action is variationally consistent
(Section~\ref{sec:G6-einstein}) and admits Schwarzschild as a
vacuum solution.  At the linearized level the propagating
graviton has the correct two transverse-traceless polarizations
(Section~\ref{sec:G4-linearized}), inherited from the gauge
structure of $h_{\mu\nu}$ as a symmetric rank-2 tensor under
diffeomorphisms of the background metric.

What the present paper does \emph{not} address:
\begin{itemize}
\item \textbf{Background independence at the path-integral level.}
      The Sakharov calculation is performed on a fixed metric;\ the
      partition function $Z = \int \mathcal{D}\psi\,
      \mathcal{D}A\,\mathcal{D}n \, e^{i S_{\mathrm{matter}}[g]}$
      treats $g$ as input, not as integration variable.  A genuinely
      background-independent NWT formulation would integrate over
      $g$ as well, and would have to demonstrate that the resulting
      effective action coincides with the Sakharov-induced one in
      a saddle-point or semi-classical limit.  We have not done
      this.
\item \textbf{Quantum-level diffeomorphism invariance.}  The
      classical effective action $\Gamma[g]$ extracted by the
      Sakharov mechanism is generally covariant by construction
      (the heat-kernel expansion organises the result in
      diffeomorphism-invariant curvature scalars).  But the
      \emph{quantization} of the emergent metric $h_{\mu\nu}$
      requires either a gauge-fixing of diffeomorphisms with
      compensating Faddeev--Popov ghosts (yielding a perturbative
      expansion that is non-renormalisable) or an independent
      argument that the substrate's projection mechanism preserves
      diffeomorphism invariance non-perturbatively.  Neither is
      supplied here.
\item \textbf{Beyond-linear-order propagating-DOF count.}  The
      two TT polarizations of the linearized graviton are
      established at $\mathcal{O}(h^2)$;\ the constraint structure
      and gauge symmetry of the full nonlinear Einstein equations
      (Section~\ref{sec:G6-einstein}) are inherited from the GR
      structure of the recovered effective action, but a direct
      derivation that the substrate's projection produces exactly
      the GR constraint algebra at all orders is not given.
\end{itemize}
These are open conceptual issues that do not affect the
leading-order classical predictions of this paper but would need to
be addressed in any quantum-gravitational treatment of the NWT
substrate.  Possible NWT-native attack routes include the
asymptotic-safety scenario for the Sakharov-induced sector, the
Volovik-style emergent-spacetime treatment of Paper~15
\cite{paper15}, and a direct quantization of the substrate algebra
itself (which would bypass the metric quantization problem by
making the substrate primary and the metric emergent).  None of
these is developed here.

% ==========================================================
\section{Conclusion}
\label{sec:conclusion}
% ==========================================================

The four-paper Null-Worldtube gravitational program extends the
structural derivation of Papers~14--17 into a dynamical theory.
Paper~14 gave the leading-order structural form
$G \propto (8/7)^2 \alpha^{21} / \me^2$ from the heptafoil Wilson
amplitude.  Paper~15 placed it on the Brieskorn--Poincar\'e sphere
$S^3/2I$ as a one-loop Casimir computation, identifying the BPS
trefoil as the relevant background.  Paper~16 supplied the
three-field Lagrangian (a complex scalar coupled to a $\U(1)$
gauge field and an $S^2$-valued Skyrme--Faddeev field at the BPS
critical coupling) on which the loop is taken, and reported the
$10^{45}$ Phase~5 structural mismatch with the experimentally
required $1/G$.  Paper~17 supplied the closed form
$\me/\mPl = (8/7)\,\alpha^{21/2}\,(1 + \alpha/7 + 3\alpha^2)$
through $K_7$ graph-state information theory, agreeing with CODATA
to $-0.0001\%$.

The present paper supplies dynamical content.  Integrating out the
BPS matter sector at one loop on an arbitrary curved background
generates the Einstein--Hilbert action via Sakharov's mechanism,
with the coefficient $1/(16\pi G)$ set by the heat-kernel $a_2$
coefficient of the matter spectrum.  Equating the resulting
Sakharov coefficient with Paper~17's structural form for $\me/\mPl$
gives the structural relation
$\Lambda_{\mathrm{UV}}/\mPl =
\sqrt{12\pi/N_{\mathrm{DOF}}^{\mathrm{sub}}}$, with
$N_{\mathrm{DOF}}^{\mathrm{sub}}$ a $c_J$-weighted sum over
\emph{substrate} matter fields (the bosonic three-field Lagrangian
of Paper~16 plus the $\so(10)$ GUT extension of Paper~14).  IR
fermions and massless gauge bosons are topological projections of
the substrate, not new fundamental fields, and do not enter the
heat-kernel sum directly;\ identifying the substrate count with one
mode per coset generator gives
$N_{\mathrm{DOF}}^{\mathrm{sub}} = 6 + 33 = 39$ and
$\Lambda_{\mathrm{UV}} = 0.983\,\mPl$, consistent with
$\Lambda_{\mathrm{UV}} \approx \mPl$ to $1.7\%$.  The Paper~17
identity $\mPl^2/\me^2 = |A_{K_7}|^{-2}$ is structurally identified
with the missing $K_7$ Lie-theoretic factor that the bare
$S^3/2I$ Casimir framework of Paper~16 Phase~5 did not include;\
a corrected Phase~5 calculation incorporating that factor is
sketched in Section~\ref{sec:G3-phase5-sketch}.

Variation of the resulting effective action $\Gamma[g]$ with
respect to $g^{\mu\nu}$ yields the full nonlinear Einstein
equations
$G_{\mu\nu} + \Lambda_{\mathrm{cc}}\,g_{\mu\nu} = 8\pi G\,T_{\mu\nu}
+ \alpha_S\,\partial_\mu\partial_\nu R + \cdots$,
with $G$ matching Paper~17's structurally derived value;\ matter
energy-momentum conservation $\nabla^\mu T_{\mu\nu} = 0$ follows
as a consistency consequence of the contracted Bianchi identity.
The Schwarzschild metric is recovered as a vacuum solution by
direct symbolic computation (all $10$ independent components of
the Ricci tensor vanishing identically).  Linearisation around
flat space reproduces the gravitational-wave equation
$\Box\,h^{\mathrm{TT}}_{\mu\nu} = -16\pi G\,T^{\mathrm{TT}}_{\mu\nu}$
with two physical transverse-traceless polarizations, and the
static weak-field limit recovers Newton's law with the
structurally derived $G$.  Numerically, the resulting prediction
matches CODATA $G$ to $+29$\,ppm, inside the $\pm 22$\,ppm
experimental error bar;\ this numerical agreement is a
re-expression of Paper~17's $\me/\mPl$ closed form via
$G\,\me^2 = \mathrm{bracket}^2\,\alpha^{21}$, not an independent
re-derivation of $G$ at the dynamical level.

The structural picture left at the end of this paper is that the
gravitational coupling $G$ first established structurally in
Papers~14--17 admits a dynamical effective action consistent with
the same number, and that this effective action coincides with
Einstein gravity at the long-wavelength limit.  The $K_7$
amplitude that fixed $\me/\mPl$ at the structural level is the
same amplitude that, together with the Sakharov coefficient,
specifies the cutoff in IR-measurable terms;\ this is a single
algebraic identity read in two complementary ways, not two
independent computations of the same quantity.  Gravity in NWT is
therefore identified, at the level developed in this paper, as the
long-wavelength variational dynamics of the Sakharov-induced
effective action of the substrate matter sector.

Four open problems remain (Section~\ref{sec:open}): the
$\Lambda_{\mathrm{cc}}$ fine-tuning, dynamical black-hole formation
from NWT vortex collapse, the higher-curvature ghost-freedom of the
Stelle sector beyond linear order, and a quantum-level treatment of
diffeomorphism invariance and background independence on curved
backgrounds.  None are NWT-specific.  The first two motivate
Papers~21 and~22 of the program; the third is a generic issue in
higher-derivative gravity that any specific realisation, including
ours, inherits; the fourth is the standard quantum-gravity problem
that any classical-limit derivation leaves open.

The next steps for the gravitational program are observational.
The companion paper~\cite{paper19} catalogues the lab- and
cosmology-scale tests in which the Sakharov-induced $G$ derived
here participates as one entry of the substrate-monism master
table; the FLRW null-test and substrate-anisotropy
sidereal-modulation programs there are the most immediate
falsifiability targets.

% ==========================================================
\section*{Acknowledgements}
% ==========================================================

We acknowledge the use of large language models (ChatGPT, Gemini,
Perplexity, and Claude across mobile and CLI surfaces) for
critical review during manuscript preparation of this paper and
earlier papers in the series.

% ==========================================================
\section*{Code availability}
% ==========================================================

The complete \texttt{nwt-substrate} library and all paper-supporting
reproduction scripts are available at
\url{https://github.com/JimGalasyn/nwt-substrate}, archived on
Zenodo~\cite{nwt_substrate_repo}.  This paper cites version
\texttt{v0.1.1}
(DOI~\href{https://doi.org/10.5281/zenodo.20012121}{10.5281/zenodo.20012121});
the concept DOI \href{https://doi.org/10.5281/zenodo.20012027}{10.5281/zenodo.20012027}
resolves to the latest release.

Specific paper-supporting components:
\begin{itemize}
\item Step-by-step Sakharov-derivation reproduction scripts
  (\texttt{analysis/}): \\
  \texttt{nwt\_\allowbreak paper18\_\allowbreak G1\_\allowbreak setup.py} (matter action, $T_{\mu\nu}$), \\
  \texttt{nwt\_\allowbreak paper18\_\allowbreak G2\_\allowbreak matter\_\allowbreak loop.py} (heat-kernel, $1/G$), \\
  \texttt{nwt\_\allowbreak paper18\_\allowbreak G3\_\allowbreak K7\_\allowbreak sakharov.py} ($K_7$ UV/IR bridge), \\
  \texttt{nwt\_\allowbreak paper18\_\allowbreak G4\_\allowbreak linearized\_\allowbreak einstein.py} (linearised Einstein, GW), \\
  \texttt{nwt\_\allowbreak paper18\_\allowbreak G5\_\allowbreak curved\_\allowbreak sakharov.py} (curved-background EH), \\
  \texttt{nwt\_\allowbreak paper18\_\allowbreak G6\_\allowbreak einstein\_\allowbreak variation.py} (full nonlinear Einstein, Schwarzschild).
\item Library gravity module \texttt{nwt\_substrate.gravity}
  packaging the validated result for re-use, including the symbolic
  Schwarzschild proof
  \texttt{verify\_\allowbreak schwarzschild\_\allowbreak vacuum\_\allowbreak symbolic()}
  and the CODATA comparison
  \texttt{G\_\allowbreak substrate\_\allowbreak SI()} that returns
  $6.674494\times10^{-11}$\,m$^3$\,kg$^{-1}$\,s$^{-2}$ ($+29$\,ppm of CODATA).
\end{itemize}

% ==========================================================
\begin{thebibliography}{99}
% ==========================================================

\bibitem{sakharov1968}
  A.\,D.~Sakharov.
  \emph{Vacuum quantum fluctuations in curved space and the theory
  of gravitation.}
  Sov. Phys. Dokl. \textbf{12}, 1040 (1968).

\bibitem{volovik2003}
  G.\,E.~Volovik.
  \emph{The Universe in a Helium Droplet.}
  Oxford University Press (2003).

\bibitem{visser2002}
  M.~Visser.
  \emph{Sakharov's induced gravity: a modern perspective.}
  Mod. Phys. Lett. A \textbf{17}, 977 (2002).

\bibitem{stelle1977}
  K.\,S.~Stelle.
  \emph{Renormalization of higher-derivative quantum gravity.}
  Phys. Rev. D \textbf{16}, 953 (1977).

\bibitem{vassilevich2003}
  D.\,V.~Vassilevich.
  \emph{Heat kernel expansion: user's manual.}
  Phys. Rep. \textbf{388}, 279 (2003).

\bibitem{appelquist1975}
  T.~Appelquist and J.~Carazzone.
  \emph{Infrared singularities and massive fields.}
  Phys. Rev. D \textbf{11}, 2856 (1975).

\bibitem{adler1982}
  S.\,L.~Adler.
  \emph{Einstein gravity as a symmetry-breaking effect in quantum
  field theory.}
  Rev. Mod. Phys. \textbf{54}, 729 (1982); Erratum
  \textbf{55}, 837 (1983).

\bibitem{paper14}
  J.~Galasyn and Th\'eodore.
  \emph{The integers are the output, not the input: $G$ from
  heptafoil topology (Paper 14).}
  Zenodo (2026), \href{https://doi.org/10.5281/zenodo.19654507}
  {10.5281/zenodo.19654507}.

\bibitem{paper15}
  J.~Galasyn and Th\'eodore.
  \emph{One knot, all forces: Casimir-induced gravity from the
  trefoil on $S^3/2I$ (Paper 15).}
  Zenodo (2026), \href{https://doi.org/10.5281/zenodo.19701263}
  {10.5281/zenodo.19701263}.

\bibitem{paper16}
  J.~Galasyn and Th\'eodore.
  \emph{The NWT Lagrangian: gravitational bootstrap from three knots
  (Paper 16).}
  Zenodo (2026), \href{https://doi.org/10.5281/zenodo.19710846}
  {10.5281/zenodo.19710846}.

\bibitem{paper17}
  J.~Galasyn and Th\'eodore.
  \emph{Newton's constant and rest-frame Schr\"odinger evolution
  from $K_7$ graph-state information theory (Paper 17).}
  Zenodo submission prepared (2026).

\bibitem{paper19}
  J.~Galasyn and Th\'eodore.
  \emph{Substrate monism: one algebra, seven vocabularies,
  all forces (Paper 19).}
  Zenodo (2026), in preparation; companion paper to the present
  one.

\bibitem{nwt_substrate_repo}
  J.~Galasyn and Th\'eodore.
  \emph{\texttt{nwt-substrate}: a substrate-algebraic computation
  library for Null Worldtube Theory (version 0.1.1).}
  Zenodo (2026).
  DOI:~\href{https://doi.org/10.5281/zenodo.20012121}{10.5281/zenodo.20012121}
  (this version);
  \href{https://doi.org/10.5281/zenodo.20012027}{10.5281/zenodo.20012027}
  (concept DOI, all versions).
  Source:~\url{https://github.com/JimGalasyn/nwt-substrate}.

\end{thebibliography}

\end{document}
